\documentclass[11pt,oneside,openany]{book}
\usepackage[T1]{fontenc}
\usepackage{lmodern,microtype}
\usepackage[a4paper,margin=28mm]{geometry}
\usepackage{amsmath,amssymb,amsthm,mathtools}
\usepackage[hidelinks,bookmarksopen,bookmarksnumbered]{hyperref}
\hypersetup{pdftitle={Finite Combinatorics: Enumeration, Graphs, and Matroids},
  pdfauthor={MathGenome},pdfsubject={Unified mathematical notes, P2-04}}
\allowdisplaybreaks[1]
\setlength{\emergencystretch}{2em}
\newcommand{\bracketbinom}[2]{\genfrac{[}{]}{0pt}{}{#1}{#2}}
\newcommand{\bracebinom}[2]{\genfrac{\{}{\}}{0pt}{}{#1}{#2}}

\theoremstyle{plain}
\newtheorem{theorem}{Theorem}[chapter]
\newtheorem{lemma}[theorem]{Lemma}
\newtheorem{proposition}[theorem]{Proposition}
\newtheorem{corollary}[theorem]{Corollary}
\theoremstyle{definition}
\newtheorem{definition}[theorem]{Definition}
\newtheorem{example}[theorem]{Example}
\theoremstyle{remark}
\newtheorem{remark}[theorem]{Remark}

\begin{document}
\frontmatter
\hypersetup{pageanchor=false}
\begin{titlepage}
\thispagestyle{empty}
\vspace*{27mm}
{\large MathGenome \enspace / \enspace P2-04\par}
\vspace{8mm}
\hrule height 0.8pt
\vspace{11mm}
{\Huge\bfseries Finite Combinatorics\par}
\vspace{4mm}
{\LARGE Enumeration, Graphs, and Matroids\par}
\vspace{12mm}
{\large Unified mathematical notes\par}
\vspace{6mm}
{\normalsize Finite counting, formal series, and symmetry\\
Graph structure, flows, matchings, and colorings\\
Matroid axioms, representation, and greedy optimization\par}
\vfill
{\large 8 November 2026\par}
\end{titlepage}
\hypersetup{pageanchor=true}
\tableofcontents

\chapter*{Scope and conventions}
\addcontentsline{toc}{chapter}{Scope and conventions}
The first five chapters develop counting, incidence algebras, formal
generating functions, partitions, and orbit enumeration. The graph
chapters give structural, enumerative, and optimization results. The
last three chapters develop finite matroids and connect their
representations and rank functions to graphs, matchings, and greedy
optimization. Stanley, Bondy--Murty, and Oxley guide the three strands
\cite{Stanley,BM,Oxley}. The result map records the eighty writing groups,
their source contracts, and the corresponding statement and proof locations.

A sum or product over an empty family has its usual neutral value.
A formal series is interpreted coefficientwise. Its coefficient ring
and the conditions for division or composition are stated when used.
The graph chapters use finite simple undirected graphs unless an
edge-labelled multigraph or a directed network is specified. The
terms source and sink in a network designate two distinct vertices.
A matroid has a finite specified ground set. Bases maximize
cardinality; weights are introduced separately in the last chapter.

The common prerequisites are finite-dimensional linear algebra and
elementary group theory. The linear-algebra interface is recorded in
\cite{D05}, and the orbit--stabilizer count is also proved in
\cite{D07}. The chapters give the determinant and representation
arguments at the points where their precise forms are needed.
The real-capacity flow proof establishes attainment by a
coordinatewise subsequence argument. The plane graph
chapter identifies the exact topology used from \cite{P202}. The divisor
specialization of incidence inversion is compared with the independent
arithmetic proof in \cite{P203}. Other textbook references give
attribution or parallel treatments.

\mainmatter
\chapter[Counting and poset inversion]{Finite counting, partially ordered sets, and inversion}
\label{p204:ch01}

We use finite sets throughout this chapter. The notation \(0!=1\)
corresponds to the unique ordering of the empty set. An empty sum is zero
and an empty product is one. When a counting identity is used in a
coefficient ring, its integer coefficients are mapped into that ring.

\section{Bijections and counting in two ways}

\begin{proposition}[Cardinality and maps]\label{p204:W01}
For finite sets \(A,B\), a bijection implies \(|A|=|B|\), an
injection \(A\to B\) implies \(|A|\le |B|\), and a surjection
implies \(|B|\le |A|\). If \(|A|=|B|\), a given map \(A\to B\)
is injective if and only if it is surjective, equivalently bijective.
\end{proposition}
\begin{proof}
A bijection transports any enumeration of \(A\) to one of \(B\).
An injection gives a bijection onto its image, and the image and its
complement partition \(B\). For a surjection, its nonempty fibers
partition \(A\); adding their sizes shows that \(|A|\ge |B|\).
If the sizes are equal and the map is injective, the complement of
the image has size zero. If the map is surjective, equality in the
fiber count forces every fiber to have size one. The empty sets obey
these same image and fiber conventions.
\end{proof}

\begin{theorem}[Finite inclusion--exclusion]\label{p204:W02}
Let \((A_i)_{i\in I}\) be a family of finite sets with finite index set.
Then
\[
 \left|\bigcup_{i\in I} A_i\right|
 =\sum_{\varnothing\ne J\subseteq I}
       (-1)^{|J|+1}\left|\bigcap_{j\in J} A_j\right|.
\]
The formula includes \(I=\varnothing\).
\end{theorem}
\begin{proof}
Sum the right side by elements of the union. If an element occurs
in the sets with indices in \(K\), it contributes the sum of
\((-1)^{|J|+1}\) over nonempty subsets of \(K\). The alternating
sum over all subsets of a nonempty \(K\) is zero, since toggling any
fixed index reverses the sign. Thus the nonempty-subset sum with
the displayed signs is one. For the empty family both sides are
zero.
\end{proof}

\begin{proposition}[Double counting an incidence relation]\label{p204:W03}
For \(R\subseteq A\times B\),
\[
 |R|=\sum_{a\in A}|\{b:(a,b)\in R\}|
     =\sum_{b\in B}|\{a:(a,b)\in R\}|.
\]
\end{proposition}
\begin{proof}
Each sum is the cardinality of a disjoint union of fibers, with
the fiber tagged by the relevant coordinate. The map to \(R\)
inserts the tag as that coordinate. Its inverse takes the coordinate
and the remaining entry from the pair.
\end{proof}

\begin{theorem}[Binomial and multinomial counting, and derangements]
\label{p204:W04}
The number of \(k\)-subsets of an \(n\)-element set is
\(\binom nk\), with value zero for \(k>n\). If \(k_1+\cdots+k_r=n\),
the number of ordered blocks of these sizes partitioning the set is
\(n!/(k_1!\cdots k_r!)\). Consequently, for commuting \(x_i\),
\[
 (x_1+\cdots+x_r)^n
 =\sum_{k_1+\cdots+k_r=n}
       \frac{n!}{k_1!\cdots k_r!}x_1^{k_1}\cdots x_r^{k_r}.
\]
The case \(r=2\) is the binomial theorem, and the binomial coefficients
satisfy Pascal's recurrence. The number \(d_n\) of permutations of
\([n]=\{1,\ldots,n\}\) without fixed points is
\[
 d_n=\sum_{j=0}^n(-1)^j\binom nj(n-j)!
     =n!\sum_{j=0}^n\frac{(-1)^j}{j!}.
\]
\end{theorem}
\begin{proof}
There are \(n!\) enumerations, since the numbers of choices at
successive positions are \(n,n-1,\ldots,1\). To produce ordered blocks,
cut an enumeration at the prescribed sizes and forget the order inside
each block. Each ordered partition has \(k_1!\cdots k_r!\) preimages.
In particular \(r=2\) gives \(k\)-subsets. For \(n\ge1\), distinguishing
whether a fixed element belongs to the subset gives
\(\binom nk=\binom{n-1}{k}+\binom{n-1}{k-1}\), with entries outside
\(0\le k\le n-1\) understood as zero. Each term of the product of
\(n\) copies of \(x_1+\cdots+x_r\) chooses a variable at each position.
The positions assigned to each variable are precisely the ordered
blocks. For the derangement count, let the forbidden event indexed by
\(i\) be fixing \(i\). Fixing all of a \(j\)-element subset leaves
\((n-j)!\) permutations. Subtract the union of the forbidden events
from all permutations, using Theorem~\ref{p204:W02}. The formula at
\(n=0\) counts the empty permutation.
\end{proof}

\begin{theorem}[The two Stirling arrays]\label{p204:W05}
Write \(\bracketbinom{n}{k}\) for the number of permutations of \([n]\) with
\(k\) cycles and \(\bracebinom{n}{k}\) for the number of partitions of \([n]\)
into \(k\) nonempty blocks. The entry at \((0,0)\) is one in both
arrays, and an out-of-range entry is zero. The recurrences are
\[
 \bracketbinom{n+1}{k}=\bracketbinom{n}{k-1}+n\bracketbinom{n}{k},
 \qquad
 \bracebinom{n+1}{k}=\bracebinom{n}{k-1}+k\bracebinom{n}{k}.
\]
For \(x^{\underline k}=x(x-1)\cdots(x-k+1)\), with
\(x^{\underline0}=1\),
\[
 x^{\underline n}=\sum_{k=0}^n(-1)^{n-k}\bracketbinom{n}{k}x^k,
 \qquad
 x^n=\sum_{k=0}^n \bracebinom{n}{k}x^{\underline k}.
\]
Thus the corresponding two lower-triangular integer arrays are
mutual inverses.
\end{theorem}
\begin{proof}
In a permutation, \(n+1\) is a singleton cycle or it is inserted
after one of the \(n\) existing elements in its cycle. Deletion of
\(n+1\) recovers the old permutation and the chosen predecessor.
In a set partition, it is a singleton block or lies in one of the
\(k\) blocks of the old partition. These are the recurrences.

For the first polynomial identity multiply the identity at \(n\) by
\(x-n\). The coefficient of \(x^k\) becomes
\[
 (-1)^{n+1-k}\bigl(\bracketbinom{n}{k-1}+n\bracketbinom{n}{k}\bigr).
\]
The initial polynomial is one. For the second identity, specialize
\(x\) to any nonnegative integer \(m\) and count maps
\([n]\to[m]\). A map with \(k\) nonempty fibers first partitions its
domain into \(k\) fibers, then assigns distinct values to those blocks
in \(m^{\underline k}\) ways. This proves equality at infinitely
many integers. The polynomial difference must be zero: over
\(\mathbb Q\), a nonzero polynomial has at most its degree many
roots, since each root gives a linear factor by polynomial division.
For each finite degree the descending factorials are a basis,
by their leading coefficients and distinct degrees. The two
identities express inverse changes of basis.
\end{proof}

\section{Lattices and incidence functions}

\begin{theorem}[Finite distributive lattices]\label{p204:W06}
Let \(L\) be a finite distributive lattice with bottom element \(0\).
Let \(J\) be the nonzero join-irreducible elements, ordered as in \(L\).
Then \(L\) is isomorphic as a lattice to the lower sets of \(J\),
by \(x\mapsto\{j\in J:j\le x\}\).
\end{theorem}
\begin{proof}
An element \(j\ne0\) is join-irreducible if \(j=u\vee v\) forces
\(j=u\) or \(j=v\). By induction on the size of its lower set, each
\(x\) is a join of join-irreducibles below it. For \(x=0\) the
join is empty. If \(x\ne0\) is reducible, express it as the join of
two strictly smaller elements and use induction on them.
If \(j\in J\) and \(j\le u\vee v\), distributivity gives
\(j=(j\wedge u)\vee(j\wedge v)\), so \(j\le u\) or \(j\le v\).
Induction extends this assertion to finite joins. The empty join
contains no nonzero element below it. Therefore, for any lower set
\(S\subseteq J\), the join of \(S\) has exactly the elements of \(S\)
below it among \(J\). Indeed a \(j\) below the join is below a member
of \(S\), and conversely every member of \(S\) lies below the join.
Taking joins of all join-irreducibles below an element gives the
other inverse. The maps preserve order, hence also joins and meets
by their universal order characterizations.
\end{proof}

\begin{definition}[Incidence algebra]\label{p204:W07}
For a finite poset \(P\) and a commutative ring \(R\), an incidence
function assigns a value in \(R\) to every pair \(x\le z\).
Set
\begin{gather*}
 (f*g)(x,z)=\sum_{x\le y\le z}f(x,y)g(y,z),
 \\
 \delta(x,z)=
 \begin{cases}1 & x=z,\\0& x<z,\end{cases}
 \qquad
 \zeta(x,z)=1.
\end{gather*}
With pointwise addition and this product the incidence functions form
an associative algebra with unit \(\delta\).
\end{definition}
\begin{proof}
Both bracketings of \(f*g*h\) sum
\(f(x,u)g(u,v)h(v,z)\) over \(x\le u\le v\le z\).
Finiteness permits changing the order of summation. In a product
with \(\delta\) the sole nonzero summand has \(y=x\) on the left or
\(y=z\) on the right. Bilinearity follows by distributing within
the finite sums.
\end{proof}

\begin{theorem}[Möbius inversion in both directions]\label{p204:W08}
The function \(\zeta\) has a two-sided convolution inverse \(\mu\).
It is specified recursively by
\[
 \mu(x,x)=1,\qquad
 \mu(x,z)=-\sum_{x\le y<z}\mu(x,y)\quad(x<z).
\]
For functions \(f,g:P\to R\), the following are equivalences:
\[
\begin{aligned}
 g(x)=\sum_{y\le x}f(y)
 &\quad\Longleftrightarrow\quad
 f(x)=\sum_{y\le x}\mu(y,x)g(y),\\
 g(x)=\sum_{y\ge x}f(y)
 &\quad\Longleftrightarrow\quad
 f(x)=\sum_{y\ge x}\mu(x,y)g(y).
\end{aligned}
\]
\end{theorem}
\begin{proof}
Write \(\zeta=\delta+h\). The \(k\)-fold convolution power \(h^{*k}\)
vanishes unless there is a strict chain of \(k\) steps between the two
endpoints, as follows by expanding the sum. In particular
\(h^{*k}=0\) for \(k>|P|\). The finite geometric sum
\(\sum_{k=0}^{|P|}(-1)^kh^{*k}\), with \(h^{*0}=\delta\), is a
two-sided inverse by telescoping. Its product with \(\zeta\) on the
right gives the displayed recurrence, which uniquely determines it
by induction on interval size. The integer-coefficient inverse
therefore specializes to any coefficient ring \(R\).

For the lower formula, substitute the sum for \(g(y)\) into
the proposed inverse sum for \(f(x)\). The coefficient of \(f(u)\) is
\[
 \sum_{u\le y\le x}\mu(y,x)=(\zeta*\mu)(u,x)=\delta(u,x).
\]
If instead one starts with the proposed formula for \(f\), summing
it for all \(y\le x\) produces \(\mu*\zeta\) and recovers \(g(x)\).
For the upper formula, substitution into
\(\sum_{y\ge x}\mu(x,y)g(y)\) produces \(\mu*\zeta\), and summing the
proposed formula for \(f\) upward produces \(\zeta*\mu\).
\end{proof}

\begin{proposition}[Products, Boolean intervals, and divisibility]
\label{p204:W09}
On a product of finite posets, the Möbius value is the product of
the two Möbius values. On a subset interval \(S\subseteq T\) it is
\((-1)^{|T\setminus S|}\). On the positive divisors of a positive
integer \(N\), ordered by divisibility, it is
\(\mu_{\mathrm{arith}}(b/a)\) at \(a\mid b\). Consequently the usual
arithmetic divisor inversion follows from Theorem~\ref{p204:W08}.
\end{proposition}
\begin{proof}
In a product interval the sum over intermediate points factors into
two sums. The product of the two Möbius functions is therefore the
inverse of the product zeta function, by the unit identities in the
two factors. The subset interval is a product of \(|T\setminus S|\)
two-element chains, each contributing \(-1\) between distinct
endpoints. This includes the empty product when \(S=T\).
For a finite chain, \(\mu(x,z)\) is one at equality, minus one across
one step and zero across two or more steps, by the recurrence in
Theorem~\ref{p204:W08}. By prime factorization, an interval
\(a\mid b\) of divisors is a product of such chains of lengths equal
to the exponents in \(b/a\). Thus the product is zero if one exponent
exceeds one, and otherwise is the sign determined by the number of
prime factors of \(b/a\). This is \(\mu_{\mathrm{arith}}(b/a)\).
Unique factorization of positive integers and the arithmetic
Möbius formula are proved in \cite[Chapters 1 and 3]{P203}.
Applying the lower inversion at \(N\) yields
\(f(N)=\sum_{d\mid N}\mu_{\mathrm{arith}}(N/d)g(d)\) from
\(g(N)=\sum_{d\mid N}f(d)\).
\end{proof}

For comparison with the standard enumeration terminology, see
\cite[Sections 1.1--1.3, 1.7--1.9, 2.1, and 3.3--3.8]{Stanley}.
\chapter[Formal series and recurrences]{Formal power series and linear recurrences}
\label{p204:ch02}

\section{Coefficient algebra}

\begin{definition}[Ordinary and exponential generating series]
\label{p204:W10}
Over a commutative ring \(R\) with identity, a formal power series
is a sequence of coefficients, written \(\sum_{n\ge0}a_nz^n\).
Write \([z^n]A=a_n\). The ordinary generating function of
\((a_n)\) is this formal series. When \(R\) is a commutative
\(\mathbb Q\)-algebra, in particular a field of characteristic zero,
its exponential generating function is \(\sum_{n\ge0}a_nz^n/n!\).
Equality means equality of coefficients.
\end{definition}

\begin{proposition}[Products and differentiation]\label{p204:W11}
The Cauchy product and formal derivative satisfy
\[
 [z^n](AB)=\sum_{i=0}^n a_i b_{n-i},\qquad
 [z^n]A'=(n+1)a_{n+1},\qquad (AB)'=A'B+AB'.
\]
Over a \(\mathbb Q\)-algebra the exponential series for two
sequences multiply to give the exponential series for
\(c_n=\sum_{i=0}^n\binom ni a_i b_{n-i}\).
\end{proposition}
\begin{proof}
The coefficient of degree \(n\) in the distributive product uses
exactly pairs of exponents totaling \(n\). Associativity follows
because a triple product sums over all triples of nonnegative
exponents totaling \(n\). Differentiation is defined on monomials
by \((z^k)'=kz^{k-1}\), with derivative zero at \(k=0\).
In the coefficient of degree \(n\) of the derivative of \(AB\),
each pair of exponents totaling \(n+1\) receives the factor \(n+1\).
In \(A'B+AB'\) the same pair receives the sum of its two exponents.
For exponential series, multiply the Cauchy coefficient by \(n!\);
the factor is \(n!/(i!(n-i)!)=\binom ni\).
\end{proof}

\begin{theorem}[Substitution and the two kinds of inverse]
\label{p204:W12}
If \(B(0)=0\), the composition \(A(B)=\sum_{j\ge0}a_jB^j\)
is a formal series. A series \(A\) is multiplicatively invertible
if and only if \(A(0)\) is a unit. Among zero-constant series,
\(B\) is compositionally invertible if and only if \([z]B\)
is a unit.
\end{theorem}
\begin{proof}
Since \(B\in zR[[z]]\), its \(j\)-th power belongs to
\(z^jR[[z]]\), so only \(j\le n\) contributes to degree \(n\).
If \(C(0)=0\) as well, the identity
\((A(B))(C)=A(B(C))\) follows by taking both sides modulo
\(z^{n+1}\), where all substitutions reduce to polynomials
and preserve the ideal \((z^{n+1})\). The same truncation
argument proves distributivity and the formal chain rule.
For the latter, to determine the derivative through degree \(n\)
truncate the original series through degree \(n+1\), apply the
polynomial chain rule, and compare coefficients.

The constant coefficient of \(AA^{-1}\) must be one. If \(a_0\)
is a unit, set \(c_0=a_0^{-1}\) and define
\[
 c_n=-a_0^{-1}\sum_{i=1}^n a_i c_{n-i}\quad(n\ge1).
\]
This gives \(AC=1\), and commutativity gives \(CA=1\).
For the compositional assertion, write
\(B=b_1z+b_2z^2+\cdots\). A compositional inverse with zero constant
must have linear coefficient \(b_1^{-1}\).
Assume \(b_1\) is a unit. In the equation \(B(C)=z\), the
coefficient of \(z^n\) is \(b_1 c_n\) plus a polynomial in
\(c_1,\ldots,c_{n-1}\). Thus there is a unique zero-constant
right inverse \(C\). Its linear term is a unit, so the same
construction gives \(C(D)=z\). Associativity yields
\(D=(B(C))(D)=B(C(D))=B\), showing that \(C\) is a left
inverse also.
\end{proof}

\section{Recurrences and rationality}

\begin{proposition}[Existence and uniqueness from initial data]
\label{p204:W13}
Fix \(r\ge1\), \(c_0,\ldots,c_{r-1}\in R\), and \(r\) initial
values \(a_0,\ldots,a_{r-1}\in R\). There is exactly one sequence
satisfying
\[
 a_{n+r}=\sum_{j=0}^{r-1}c_j a_{n+j}\quad(n\ge0)
\]
with these initial values.
\end{proposition}
\begin{proof}
At step \(n\), the right side uses only entries with index
less than \(n+r\). It defines the new entry from previously
defined ones. Two sequences with the same first \(r\) entries
then agree at every subsequent index by induction.
\end{proof}

\begin{proposition}[The rational ordinary generating function]
\label{p204:W14}
For a sequence as in Proposition~\ref{p204:W13}, let
\[
 Q(z)=1-\sum_{j=0}^{r-1}c_jz^{r-j},\qquad A(z)=\sum_{n\ge0}a_nz^n.
\]
Then \(QA=P\) where \(P\in R[z]\) has degree less than \(r\)
and is determined by the \(r\) initial values.
In particular \(A=P/Q\), with \(Q(0)=1\).
\end{proposition}
\begin{proof}
The coefficient of \(z^{n+r}\) in \(QA\) is the left side
of the recurrence minus its right side. It is zero for \(n\ge0\).
The coefficients of degrees \(0,\ldots,r-1\) in \(QA\) use only
the initial entries, determining \(P\). Theorem~\ref{p204:W12}
makes division by \(Q\) well defined.
\end{proof}

\begin{theorem}[Distinct characteristic roots]\label{p204:W15}
Let \(R=K\) be a field and suppose that
\[
 \chi(t)=t^r-\sum_{j=0}^{r-1}c_jt^j
\]
has \(r\) distinct roots \(\lambda_1,\ldots,\lambda_r\) in \(K\).
Every solution in Proposition~\ref{p204:W13} is a unique linear
combination of \((\lambda_i^n)_{n\ge0}\), with \(\lambda_i^0=1\)
also when \(\lambda_i=0\).
\end{theorem}
\begin{proof}
For each root, multiplying \(\chi(\lambda_i)=0\) by
\(\lambda_i^n\) verifies the recurrence, including \(n=0\)
when the root is zero. On the first \(r\) entries the \(r\)
sequences form the columns of the Vandermonde matrix with rows
indexed by exponents \(0,\ldots,r-1\).
Evaluation of polynomials of degree below \(r\) at the roots is
the transpose map. It is onto, since for each \(i\) the polynomial
\(\prod_{j\ne i}(t-\lambda_j)/(\lambda_i-\lambda_j)\) has
values one at \(i\) and zero at all other roots. The matrix is
square and hence invertible. Each set of initial values is matched
uniquely, and Proposition~\ref{p204:W13} extends equality to the
full sequence.
\end{proof}

\begin{theorem}[Eventual recurrence and roots of arbitrary multiplicity]
\label{p204:W16}
Over a field \(K\), an ordinary generating series is a quotient
of polynomials with denominator of constant term one if and
only if its sequence eventually satisfies a homogeneous
constant-coefficient recurrence. If a recurrence starts at
index \(0\) and its characteristic polynomial splits as
\(\chi(t)=\prod_\lambda (t-\lambda)^{m_\lambda}\), a basis of
the solution space is
\[
 b_{\lambda,j}(n)=[z^n]\frac{z^j}{(1-\lambda z)^{j+1}},
 \qquad 0\le j<m_\lambda .
\]
For \(n\ge j\) this value is \(\binom nj\lambda^{n-j}\), with
value zero for \(n<j\). Here the integer binomial coefficient
is mapped into \(K\). When \(\lambda=0\), the series is \(z^j\).
In characteristic zero and for \(\lambda\ne0\), the corresponding
root space consists of sequences \(P(n)\lambda^n\) with
\(\deg P<m_\lambda\).
\end{theorem}
\begin{proof}
If \(QA=P\) with \(Q(0)=1\), let \(s=\deg Q\).
For each index \(n\ge s\) above \(\deg P\), the coefficient
identity reads \(a_n=-\sum_{i=1}^s [z^i]Q\,a_{n-i}\).
If \(s=0\), the sequence is eventually zero and satisfies
the order-one recurrence with zero coefficient there.
Conversely an eventual recurrence gives a polynomial \(Q\)
as in Proposition~\ref{p204:W14}, and \(QA\) has only finitely
many potentially nonzero coefficients.

For the root formula, expanding \((1-\lambda z)^{-(j+1)}\)
as a product of geometric series gives the binomial coefficient.
Indeed the coefficient of degree \(n-j\) counts decompositions
of \(n-j\) into \(j+1\) ordered nonnegative parts, in bijection
with placements of \(j\) separators among \(n\) positions.
Let \(S\) denote the left shift \(Sa(n)=a(n+1)\).
Pascal's identity gives
\[
 (S-\lambda)b_{\lambda,0}=0,\qquad
 (S-\lambda)b_{\lambda,j}=b_{\lambda,j-1}\quad(j\ge1).
\]
This includes \(\lambda=0\), when a unit spike at \(j\) shifts
to a spike at \(j-1\), or to zero when \(j=0\).
It follows that \(\chi(S)\) kills each proposed solution.
For fixed \(\lambda\) the sequences are independent because
their first \(m_\lambda\) entries form a triangular matrix
with diagonal entries one. On their span \(S=\lambda I+N\)
where \(N\) is nilpotent. For another root \(\mu\ne\lambda\),
\((S-\mu)\) on this span is invertible, by the finite geometric
inverse of \((\lambda-\mu)I+N\).
In a linear relation among the root spans, applying
\(\prod_{\mu\ne\lambda}(S-\mu)^{m_\mu}\) kills all other
spans and is invertible on the \(\lambda\)-span. Thus the spans
are independent together. Their dimensions sum to \(r\),
which is the solution dimension by the initial-value result.
If the characteristic is zero and \(\lambda\ne0\), each
\(\lambda^{-j}\binom nj\) is a polynomial of degree \(j\) with
nonzero leading coefficient, and those polynomials form a
basis of degrees below \(m_\lambda\).
\end{proof}

\begin{proposition}[Transfer matrices]\label{p204:W17}
For a finite square matrix \(M\) over a commutative ring,
\((I-zM)^{-1}=\sum_{n\ge0}M^n z^n\).
Each entry is a rational series over a field. If \(M_{uv}\)
is the total weight of arcs \(u\to v\) in a finite directed
weighted graph, \((M^n)_{uv}\) is the sum of the products
of weights of all \(n\)-step walks from \(u\) to \(v\).
\end{proposition}
\begin{proof}
Multiplying the displayed series by \(I-zM\) on either side
cancels every positive-degree coefficient and gives \(I\)
in degree zero. The adjugate identity gives the same inverse
as \(\operatorname{adj}(I-zM)/\det(I-zM)\), whose denominator
has constant term one. The adjugate identity follows by
cofactor expansion in the diagonal entries of the matrix
product and expansion of a determinant with two equal rows
or columns in the off-diagonal entries. The walk assertion
follows from the matrix product by induction on \(n\),
partitioning walks by their penultimate vertex and last arc.
The power \(M^0=I\) represents the empty walks at the vertices.
The empty matrix case has no entries to assert.
\end{proof}

\section{Lagrange inversion}

\begin{theorem}[Lagrange--Bürmann inversion]\label{p204:W18}
Let \(R\) be a commutative \(\mathbb Q\)-algebra, and let
\(\phi,H\in R[[z]]\) with \(\phi(0)\) a unit.
There is a unique \(F\in zR[[z]]\) such that \(F=z\phi(F)\).
For \(n\ge1\),
\[
 [z^n]H(F(z))=\frac{1}{n}[u^{n-1}]H'(u)\phi(u)^n .
\]
\end{theorem}
\begin{proof}
Define \(g(u)=u/\phi(u)\). Its linear coefficient is the inverse
of \(\phi(0)\). By Theorem~\ref{p204:W12} it has a unique
compositional inverse \(F\), and \(g(F(z))=z\) is equivalent
to \(F=z\phi(F)\).

For a formal Laurent series, which has only finitely many negative
exponents, write \(\operatorname{res}_u L=[u^{-1}]L\).
For any \(g=uU(u)\) with \(U(0)\) a unit, we have
\[
 \operatorname{res}_u L(g(u))g'(u)=\operatorname{res}_z L(z).
\]
To prove this, it suffices to check monomials \(L(z)=z^k\):
only finitely many terms can contribute to the residue because
the least possible exponent of \(g^k g'\) is at least \(k\).
If \(k\ne-1\), \(g^k g'=(g^{k+1})'/(k+1)\) and a derivative
has residue zero. The power rule for negative exponents follows
from differentiating \(g\,g^{-1}=1\) and the product rule;
those rules extend to Laurent series by the same finite
coefficient computations. If \(k=-1\), then
\(g'/g=u^{-1}+U'/U\) and \(U'/U\) is an ordinary power series.
This proves the substitution identity.

Differentiate \(H(F(z))\), then express the coefficient as a residue.
The substitution identity gives
\[
\begin{aligned}
 n[z^n]H(F(z))
 &=\operatorname{res}_z z^{-n}H'(F(z))F'(z)\\
 &=\operatorname{res}_u g(u)^{-n}H'(u)F'(g(u))g'(u)\\
 &=\operatorname{res}_u u^{-n}H'(u)\phi(u)^n .
\end{aligned}
\]
The last step uses \(F'(g(u))g'(u)=1\) and \(g=u/\phi\).
Reading the coefficient and dividing by \(n\) proves the assertion.
\end{proof}

For comparison, see \cite[Chapter 7]{GKP},
\cite[Sections 4.1, 4.3, and 4.7]{Stanley}, and
\cite[Appendix A.6]{FS}.
\chapter[Combinatorial constructions]{Combinatorial constructions, Catalan numbers, and labelled sets}
\label{p204:ch03}

\section{The basic counting constructions}

An unlabelled class is specified here by finite sets \(\mathcal A_n\)
of objects of each size \(n\). A labelled class also specifies its
objects on each finite label set \(S\), with consistent transports
under bijections of label sets. We count on a fixed \(n\)-element
label set, so the number \(a_n\) is independent of which labels are
chosen. The product of two labelled classes uses a disjoint split
of the label set. See \cite[Chapters I and II]{FS} for this notation.

\begin{proposition}[Disjoint sum and product]\label{p204:W19}
For locally finite classes, a disjoint sum adds the counting sequences.
The product with additive size gives
\(\sum_{i=0}^n a_i b_{n-i}\) in the unlabelled case.
The labelled product gives
\(\sum_{i=0}^n\binom ni a_i b_{n-i}\).
Thus the ordinary series multiply for the unlabelled product, and
the exponential series multiply for the labelled product.
\end{proposition}
\begin{proof}
The two summands are tagged by their class. For an unlabelled
product of total size \(n\), record the size \(i\) of the first
object, and count the two choices independently. In the labelled
case, the first object is given an \(i\)-subset of the \(n\) labels,
and the second uses its complement. Consistent transport gives
\(a_i\) and \(b_{n-i}\) choices on the two chosen sets, respectively.
The stated series identities follow from
Proposition~\ref{p204:W11}.
\end{proof}

\begin{theorem}[Ordered sequences and labelled sets]\label{p204:W20}
Let \(\mathcal A\) be locally finite and have no size-zero objects.
The ordinary series of finite ordered sequences of unlabelled
\(\mathcal A\)-objects is \(1/(1-A(z))\), where \(A\) is their
ordinary series. If \(\mathcal A\) is labelled and \(A\) denotes
its exponential series, finite sets of \(\mathcal A\)-objects on
disjoint label blocks covering the given labels have exponential
series \(\exp(A(z))\).
\end{theorem}
\begin{proof}
In the unlabelled case, a sequence of exactly \(k\) objects has
series \(A^k\) by repeated application of
Proposition~\ref{p204:W19}. The empty sequence contributes one.
Each coefficient involves only \(k\le n\) because all sizes
are positive. The resulting geometric sum has product one with
\(1-A\).
In the labelled case, the ordered \(k\)-tuples likewise have
exponential series \(A^k\). For an unordered object with \(k\)
components, the components have disjoint nonempty label sets,
so the \(k!\) orders are all distinct and exhaust its preimages.
The term \(A^k/k!\) counts exactly the objects with \(k\)
components. Summing these coefficientwise gives the formal
exponential, including the empty set.
\end{proof}

\begin{theorem}[Subset counts and the Catalan family]\label{p204:W21}
The number of \(k\)-subsets is \(\binom nk\). A Dyck path of
semilength \(n\) is a word of \(n\) up-steps and \(n\) down-steps
whose partial height never goes below zero. These words are also
balanced-parenthesis strings when up-steps open parentheses and
down-steps close them. Their common
number \(C_n\) satisfies
\[
 C_0=1,\qquad
 C_{n+1}=\sum_{i=0}^n C_i C_{n-i},\qquad
 C_n=\frac{1}{n+1}\binom{2n}{n}.
\]
\end{theorem}
\begin{proof}
The subset count is Theorem~\ref{p204:W04}, and the
parenthesis condition is exactly the partial-height condition.
A nonempty Dyck path begins with an up-step. At the first
return to height zero, the down-step closes an interior
Dyck path, read at height one. The suffix is a second Dyck path.
This decomposition is unique and reversible, giving the
recurrence by semilength.

All balanced words without the nonnegativity restriction are
counted by \(\binom{2n}{n}\). For \(n\ge1\), reflect the prefix of
a bad word through its first visit to height \(-1\), by exchanging
up- and down-steps in that prefix. The resulting word has \(n+1\)
up-steps and \(n-1\) down-steps. Its first visit to height \(1\)
ends the reflected prefix, so reflection there gives the inverse.
Thus the bad count is \(\binom{2n}{n+1}\); subtracting and
simplifying gives the closed formula. At \(n=0\) it holds for
the unique empty word.
\end{proof}

\begin{proposition}[Labelled cycles, block counts, and Bell numbers]
\label{p204:W22}
For a locally finite labelled class with positive sizes, cycles
of components on disjoint label blocks have exponential series
\[
 \sum_{k\ge1}\frac{A(z)^k}{k}=-\log(1-A(z)).
\]
For ordinary set partitions of an \(n\)-element labelled set,
marking each block by \(u\) gives exponential series
\[
 \exp\bigl(u(e^z-1)\bigr)
 =\sum_{n\ge0}\left(\sum_{k=0}^n \bracebinom{n}{k}u^k\right)\frac{z^n}{n!}.
\]
The Bell number \(B_n=\sum_{k=0}^n\bracebinom{n}{k}\) therefore has
exponential series \(\exp(e^z-1)\).
\end{proposition}
\begin{proof}
An ordered \(k\)-tuple of labelled components has exponential
series \(A^k\). Passing to a cyclic order identifies precisely
the \(k\) rotations of the tuple. A nontrivial rotation cannot
fix it, since the label blocks are disjoint and nonempty.
Thus division by \(k\) counts cyclic orders. Finiteness of each
coefficient again follows from positive sizes; the logarithm
is its formal series expansion. For set partitions, there is
one nonempty unstructured block on each nonempty label set,
so a single block contributes \(u\sum_{n\ge1}z^n/n!\).
The labelled-set formula in Theorem~\ref{p204:W20} then
gives the series, with the coefficient of \(u^kz^n/n!\)
counting the \(k\)-block partitions. Setting \(u=1\) proves
the last claim.
\end{proof}

\section{A determinant for disjoint paths}

\begin{theorem}[Lindström--Gessel--Viennot]\label{p204:W23}
Let a finite acyclic directed graph have commuting edge weights
in a commutative ring. Choose vertices
\(a_1,\ldots,a_r\) and \(b_1,\ldots,b_r\), distinct within
each list, and let \(M_{ij}\) be the sum of path weights from
\(a_i\) to \(b_j\), where a path weight is the product of its edge
weights. Empty paths have weight one. Then
\[
 \det M
 =\sum_{\substack{\sigma\in S_r,\ (P_1,\ldots,P_r)\\
                  P_i:a_i\to b_{\sigma(i)}\\
                  P_i\text{ pairwise vertex-disjoint}}}
          \operatorname{sgn}(\sigma)\prod_{i=1}^r w(P_i).
\]
If every disjoint family must use \(\sigma=\operatorname{id}\),
the formula is the unsigned weight sum for that matching.
\end{theorem}
\begin{proof}
A finite acyclic digraph has an order of vertices in which every
edge points forward. There is a vertex with no incoming edge,
since otherwise following predecessors would eventually repeat
a vertex and give a directed cycle. Repeated removal of such
a vertex constructs the order. In particular each directed
walk is a path and has finite length bounded by the number of
vertices minus one.

Expand the determinant, then each product of path sums. In a
family with an intersection, choose the first vertex in the
fixed order lying on at least two paths, and the lexicographically
first pair of path indices through it. Switch the two tails
after the chosen vertex, keeping the prefixes up to that vertex.
Each new path is valid and has no repetition, since its prefix
uses vertices no later than the chosen vertex and its tail
uses only later vertices. The first intersection and chosen
pair are unchanged: only the suffixes at strictly later vertices
have moved. Applying the same operation again recovers the
original family. The two target indices have been transposed,
so the sign is reversed. The product of weights is preserved,
as the multiset of edges in the whole family is unchanged.
All intersecting families therefore cancel, leaving the stated
formula. For \(r=0\), the empty family contributes the empty
determinant \(1\).
\end{proof}

For the determinant construction and its enumerative use,
see \cite{LGV}.
\chapter[Partitions and Gaussian coefficients]{Integer partitions and Gaussian binomial coefficients}
\label{p204:ch04}

\section{Partitions and Euler's product}

\begin{definition}[Partitions and multiplicities]\label{p204:W24}
A partition of \(n\ge0\) is a weakly decreasing finite list
\(\lambda=(\lambda_1,\ldots,\lambda_k)\) of positive integers
with sum \(n\). Its multiplicities are
\(m_i=|\{j:\lambda_j=i\}|\), a finitely supported sequence with
\(\sum_{i\ge1} i m_i=n\).
These two presentations are bijective. The empty list is the
only partition of zero.
\end{definition}
\begin{proof}
A list determines its multiplicities. From the multiplicities,
write each positive \(i\) exactly \(m_i\) times, in decreasing
order of \(i\). This is finite and reverses the first map.
When \(n=0\), the equality \(\sum_i i m_i=0\) forces every
multiplicity to vanish. For any fixed \(n\), only \(i\le n\)
can occur and \(m_i\le n/i\), so there are finitely many partitions.
\end{proof}

\begin{theorem}[Euler's partition product]\label{p204:W25}
Write \(p(n)\) for the number of partitions of \(n\). As formal
series over \(\mathbb Z\), and hence over \(\mathbb Q\),
\[
 \sum_{n\ge0}p(n)z^n=\prod_{k\ge1}(1-z^k)^{-1}.
\]
The product is defined by stabilization of coefficients of its
finite initial products.
\end{theorem}
\begin{proof}
For the part \(k\), the factor is
\(\sum_{m\ge0}z^{km}\), since multiplication by \(1-z^k\)
gives one coefficientwise. A term of degree \(n\) in a product
of the first \(N\) factors chooses a multiplicity for each
\(k\le N\) whose total weight is \(n\). If \(N\ge n\), no
further factor has a nonconstant term in degrees at most \(n\).
The stabilized choices are exactly the multiplicity presentations
of Definition~\ref{p204:W24}.
\end{proof}

\begin{theorem}[Odd parts and distinct parts]\label{p204:W26}
For each \(n\ge0\), partitions of \(n\) into odd parts are
equinumerous with partitions of \(n\) into distinct parts.
\end{theorem}
\begin{proof}
For each odd part \(d\), expand the multiplicity \(m_d\)
in binary:
\[
 m_d=\sum_{j\ge0}\varepsilon_{d,j}2^j,\qquad
 \varepsilon_{d,j}\in\{0,1\}.
\]
Existence and uniqueness of the digits follow
by repeated division by two, since the least digit is the
remainder. Replace all copies of \(d\) by the parts \(2^jd\)
for which the digit is one. This preserves total weight.
Each positive integer has a unique decomposition as \(2^jd\)
with \(d\) odd, so the resulting parts are distinct and the
inverse groups them by their odd parts. The empty list is
fixed.
\end{proof}

\begin{proposition}[Ferrers conjugation]\label{p204:W27}
The Ferrers diagram of \(\lambda\) is the set of lattice cells
\[
 \{(i,j):1\le i\le k,\ 1\le j\le\lambda_i\}.
\]
Interchanging the two coordinates defines a weight-preserving
involution on partitions. It exchanges the number of parts with
the largest part, taking both statistics to be zero on the empty
partition.
\end{proposition}
\begin{proof}
The length of row \(j\) after transposition is the number of
original rows of length at least \(j\). These lengths are positive
for \(j\le\lambda_1\) and are weakly decreasing, so this is
a partition. Transposition is its own inverse and preserves the
number of cells. The nonempty transposed diagram has \(\lambda_1\)
rows and its first row has length \(k\). The empty case follows
directly.
\end{proof}

\section{The rectangular and signed refinements}

\begin{theorem}[Rectangular partitions and the finite \(q\)-binomial identity]
\label{p204:W28}
Let \(F(a,b)\in\mathbb Z[q]\) be the weight polynomial
\(\sum_\lambda q^{|\lambda|}\) for partitions with at most \(a\)
parts and largest part at most \(b\). For \(a,b\ge1\),
\[
 F(a,b)=F(a,b-1)+q^b F(a-1,b),\qquad
 F(a,0)=F(0,b)=1.
\]
For \(0\le k\le n\), put \(\bracketbinom{n}{k}_q=F(k,n-k)\), and set it
to zero outside this range. Then
\[
 \bracketbinom{n}{k}_q=\prod_{j=1}^k\frac{1-q^{n-k+j}}{1-q^j},
 \qquad
 \prod_{i=0}^{n-1}(1+tq^i)=
     \sum_{k=0}^n q^{k(k-1)/2}\bracketbinom{n}{k}_q t^k .
\]
The first equality identifies a quotient in \(\mathbb Q(q)\)
with a polynomial.
\end{theorem}
\begin{proof}
Either no part has size \(b\), or one may remove the first
part of size \(b\). The latter operation is inverted by adding
one such part at the start. This gives the recurrence and the
boundary values. For \(G(n,k)=\bracketbinom{n}{k}_q\) the recurrence
becomes
\[
 G(n,k)=G(n-1,k)+q^{n-k}G(n-1,k-1).
\]
To verify it for the product expression \(P(n,k)\), for
\(0<k<n\) divide \(P(n-1,k)\) and \(P(n-1,k-1)\) by \(P(n,k)\).
The ratios are respectively
\((1-q^{n-k})/(1-q^n)\) and \((1-q^k)/(1-q^n)\).
Their weighted sum in the recurrence is one. The boundary
products at \(k=0,n\) are one, so induction proves \(P=G\).
For the last formula, multiply the claimed identity at \(n-1\)
by \(1+tq^{n-1}\). After factoring \(q^{k(k-1)/2}\) in
the coefficient of \(t^k\), the second summand has factor
\(q^{n-k}\), giving exactly the recurrence. At \(n=0\)
the product and sum both equal one.
\end{proof}

\begin{theorem}[Euler's pentagonal theorem and partition recurrence]
\label{p204:W29}
In \(\mathbb Z[[z]]\),
\[
 \prod_{i\ge1}(1-z^i)
    =\sum_{j\in\mathbb Z}(-1)^j z^{j(3j-1)/2}.
\]
With \(p(0)=1\) and \(p(m)=0\) for \(m<0\), this implies for \(n>0\)
\[
 p(n)=\sum_{k\ge1}(-1)^{k+1}
 \left(p\bigl(n-k(3k-1)/2\bigr)+p\bigl(n-k(3k+1)/2\bigr)\right).
\]
\end{theorem}
\begin{proof}
Expand the finite initial products and stabilize each coefficient
as in Theorem~\ref{p204:W25}. This yields the signed weight sum
of distinct-part partitions, where the sign is \((-1)^k\)
for \(k\) parts. For a nonempty decreasing list of distinct
parts \(\lambda_1>\cdots>\lambda_k>0\), write \(s=\lambda_k\)
and let \(t\) be the length of the initial consecutive run,
so \(\lambda_i=\lambda_1-i+1\) for \(1\le i\le t\).
The following two operations are inverse except in the cases
described below.

If \(s>t\), subtract one from each of the first \(t\) parts
and append a part of size \(t\). When \(t<k\), maximality of
the run makes the gap at \(t\) at least two before subtraction,
and the new last part is strictly smaller than \(s\).
When \(t=k\), the new list is strictly decreasing if and only
if \(s-1>t\). Thus the only failure in this case is \(t=k\),
\(s=k+1\). In every successful case, the new smallest part
is \(t\) and the new initial run has length at least \(t\).

If \(s\le t\), remove the last part and add one to each of
the first \(s\) remaining parts. There are at least \(s\)
remaining parts unless \(s=t=k\), the other failure. When
successful, the initial run of the new list has length exactly
\(s\). The first \(s\) parts remain consecutive. If there
is a following part, the gap at that point increases by one,
ending the run. The new smallest part exceeds \(s\). If all
remaining parts were increased, the new smallest part is at
least \(s+2\), so the inverse subtraction and appending
operation is valid even at the end of the run.

After the first operation, the new smallest part is at most
the new run length, so the second operation removes precisely
the appended part and restores the decreased entries. After
the second, the new smallest part exceeds the new run length,
so the first restores the removed part and the previous entries.
The operations preserve total weight and change the number of
parts by one. Their contributions therefore cancel in each
coefficient.

In the failure cases the whole list is consecutive. For \(s=k\)
its weight is \(k+(k+1)+\cdots+(2k-1)=k(3k-1)/2\). For
\(s=k+1\) the weight is \(k(3k+1)/2\). Each has sign
\((-1)^k\), and the empty list contributes one. These are the
terms indexed by \(j=k\), \(j=-k\), and \(j=0\) in the identity.
Multiplying by the inverse product in Theorem~\ref{p204:W25}
gives one. Equating coefficients for \(n>0\) and isolating
the \(j=0\) term gives the recurrence, with only finitely many
nonzero summands.
\end{proof}

See \cite[Sections 1.7--1.8]{Stanley} and \cite{Postnikov} for
the historical forms of the partition and \(q\)-binomial identities.
\chapter[Group actions and Pólya counting]{Finite group actions and Pólya enumeration}
\label{p204:ch05}

\section{From fixed points to orbits}

Let \(G\) be a finite group acting on a finite set \(X\) on the left.
Write \(\operatorname{Fix}_X(g)=\{x:gx=x\}\) and
\(G_x=\{g:gx=x\}\). Orbits partition \(X\).

\begin{theorem}[Burnside's fixed-point formula]\label{p204:W30}
For any such action,
\[
 |G|\,|X/G|=\sum_{g\in G}|\operatorname{Fix}_X(g)|.
\]
There is no faithfulness requirement, and \(X\) may be empty.
\end{theorem}
\begin{proof}
Count the pairs \((g,x)\) with \(gx=x\) by their two coordinates
as in Proposition~\ref{p204:W03}. The count by \(x\) is
\(\sum_x|G_x|\). Fix an orbit \(O\) and \(x\in O\).
The map \(G\to O\), \(g\mapsto gx\), is onto. If \(gx=hx\),
then \(h^{-1}g\in G_x\), so each fiber is a left coset of
\(G_x\) and has \(|G_x|\) elements. Thus \(|O||G_x|=|G|\).
Stabilizers along the orbit have the same size by conjugation,
or by the same fiber count at each point. Summing over \(O\)
therefore contributes \(|G|\). Summing over the orbits proves the
identity. This also proves the orbit--stabilizer count used below.
\end{proof}

\begin{definition}[Cycle index, including fixed points]\label{p204:W31}
For \(g\in G\), let \(c_k(g)\) be the number of length-\(k\)
cycles of the permutation \(x\mapsto gx\) of \(X\).
The cycle index is the polynomial
\[
 Z_{G,X}(s_1,\ldots,s_{|X|})
  =\frac1{|G|}\sum_{g\in G}\prod_{k=1}^{|X|}s_k^{c_k(g)}
  \ \in\ \mathbb Q[s_1,\ldots,s_{|X|}].
\]
A fixed point is a cycle of length one. For \(X=\varnothing\)
the monomial for each \(g\) is one.
\end{definition}

\begin{proof}
Iteration of a finite permutation returns each \(x\) to itself,
because the iterates eventually repeat and the permutation is
invertible. The least positive return time is the size of its
cyclic orbit, and intersecting cyclic orbits are equal.
This partitions \(X\) into the cycles counted by \(c_k(g)\),
including length-one cycles. If \(H\) is the image of the action
homomorphism in the permutation group, all fibers \(G\to H\)
have size equal to the kernel size. Hence averaging over \(G\)
or over \(H\) gives the same polynomial.
\end{proof}

\begin{theorem}[Unweighted Pólya enumeration]\label{p204:W32}
For a finite color set \(C\), write \(q=|C|\).
The induced left action on the colorings \(f:X\to C\) is
\((g\cdot f)(x)=f(g^{-1}x)\). The number of coloring orbits is
\[
 Z_{G,X}(q,\ldots,q).
\]
In particular a fixed coloring for \(g\) has
\(q^{\sum_k c_k(g)}\) possible values over all its cycles.
\end{theorem}
\begin{proof}
The inverse in the definition ensures
\(g\cdot(h\cdot f)=(gh)\cdot f\). A coloring is fixed by \(g\)
if and only if it is constant on each cycle of \(g\).
Restriction to one color on each cycle is a bijection from
fixed colorings to functions from the cycle set to \(C\).
Burnside's formula now gives the stated average. If \(C\) is
empty and \(X\ne\varnothing\), each cycle set is nonempty and
there are no colorings. If \(X=\varnothing\), there is one
empty coloring even for \(C=\varnothing\), in agreement with
the empty-product convention.
\end{proof}

\begin{theorem}[Weighted cycle-index substitution]\label{p204:W33}
Introduce commuting indeterminates \(w_c\) for \(c\in C\).
The weight of a coloring is
\(W(f)=\prod_{x\in X}w_{f(x)}\), which is constant on its orbit.
Writing \(W(O)\) for the weight of an orbit, one has
\[
 \sum_{O\in C^X/G}W(O)
 =Z_{G,X}\left(\sum_{c\in C}w_c,\ \sum_{c\in C}w_c^2,\
                         \ldots,\ \sum_{c\in C}w_c^{|X|}\right).
\]
The right side, initially written over \(\mathbb Q\), has
nonnegative integer coefficients. It specializes to arbitrary
commuting weights in a commutative semiring via that integral
polynomial.
\end{theorem}
\begin{proof}
A group element permutes positions, so \(W\) is invariant.
Work first with the polynomial ring over \(\mathbb Z\).
The orbit--stabilizer fiber computation gives
\[
 \sum_{g\in G}\sum_{f:\,g\cdot f=f} W(f)
  =\sum_{f\in C^X}|G_f| W(f)
  =|G|\sum_{O\in C^X/G} W(O).
\]
For fixed \(g\), a \(k\)-cycle has one independently chosen
color and contributes \(w_c^k\). Distributivity yields
\(\prod_k(\sum_{c\in C}w_c^k)^{c_k(g)}\) for the inner sum.
Dividing the integer polynomial identity by \(|G|\) over
\(\mathbb Q\) proves the displayed formula. Its coefficients
are the numbers of orbits with each color multiplicity, so
they are nonnegative integers. Evaluation of that polynomial
uses only addition and multiplication of weights. In a semiring
where \(|G|\) is not a unit, this last evaluation does not
require any division. The argument also covers empty color
sets and empty position sets.
\end{proof}

\section{Necklaces and bracelets}

\begin{proposition}[The rotation and reflection cycle indices]
\label{p204:W34}
For \(n\ge1\), take the positions \(\mathbb Z/n\mathbb Z\).
The rotations give the necklace cycle index
\[
 Z_{\mathrm{rot},n}
 =\frac1n\sum_{d\mid n}\varphi(d)s_d^{\,n/d},
\]
where \(\varphi(d)\) counts residues modulo \(d\) coprime to \(d\).
Allowing reflections as well gives the bracelet cycle index
\[
 Z_{\mathrm{dih},n}
 =\frac12 Z_{\mathrm{rot},n}
  +\begin{cases}
     \frac12s_1s_2^{(n-1)/2},& n\ \text{odd},\\[2pt]
     \frac14\left(s_1^2s_2^{(n-2)/2}+s_2^{n/2}\right),
       & n\ \text{even}.
   \end{cases}
\]
Factors with exponent zero are omitted, including in the \(n=1\)
formula. Substitution as in Theorem~\ref{p204:W32} or
Theorem~\ref{p204:W33} gives the unweighted counts or the weight
polynomials for these orbits. At \(n=0\) we use the empty-position
action and there is one empty object.
\end{proposition}
\begin{proof}
Set \(d_0=\gcd(n,a)\) for the rotation \(x\mapsto x+a\).
Its least positive return time is \(n/d_0\). Indeed the equation
\(n\mid ma\) reduces to \(n/d_0\mid m\) after cancelling the
coprime factor \(a/d_0\). This return time \(d\) is the same
for every position, so the monomial is \(s_d^{n/d}\).
There are \(\varphi(d)\) choices of \(a\) with return time \(d\),
since \(a\) is \(n/d\) times a residue coprime to \(d\) modulo \(d\).
This includes the identity at \(d=1\). The rotation formula follows.

Use the abstract dihedral group of pairs
\((a,\varepsilon)\), \(a\in\mathbb Z/n\mathbb Z\),
\(\varepsilon\in\{1,-1\}\), with multiplication
\((a,\varepsilon)(b,\eta)=(a+\varepsilon b,\varepsilon\eta)\).
The action is \(x\mapsto a+\varepsilon x\). It is allowed to have
a kernel for small \(n\); by Definition~\ref{p204:W31} this does
not change its cycle index. The \(n\) elements with
\(\varepsilon=-1\) act as involutions. Their fixed points solve
\(2x=a\) modulo \(n\). If \(n\) is odd, multiplication by two
is a bijection, so there is one fixed point and the rest form
2-cycles. If \(n\) is even, this map has kernel of size two and
image the even residues. Half the reflections have two fixed
points, and half have none. In each case the nonfixed points
are paired. Averaging these monomials with the rotation half of
the group gives the formula. For \(n=1,2\) it still averages
the abstract group, including repetitions of the same permutation.
\end{proof}

For the group-theoretic enumeration background and applications,
see \cite{Polya}.
\chapter{Finite graphs, connectivity, and spanning trees}
\label{p204:ch06}

Graphs in this chapter are finite and simple. Thus \(G=(V,E)\) consists of
a finite vertex set \(V\) and a set \(E\) of two-element subsets of \(V\).
Write \(uv\in E\) for \(\{u,v\}\in E\), and \(d_G(v)\) for the number of
neighbors of \(v\). When a later chapter uses parallel directed arcs,
those arcs are individually indexed.

\section{Paths and trees}

\begin{definition}[Subgraphs, paths, and components]\label{p204:W35}
A subgraph has a subset of the vertices and a subset of the edges with
both ends in that subset. The induced subgraph \(G[S]\) on \(S\subseteq V\)
keeps every such edge. A walk of length \(k\geq 0\) is a sequence
\(v_0,\ldots,v_k\) with \(v_{i-1}v_i\in E\) for \(1\leq i\leq k\).
A path is a walk with all vertices distinct. A cycle is a closed walk
\(v_0,\ldots,v_k=v_0\) with \(k\geq 3\) and \(v_0,\ldots,v_{k-1}\)
distinct. A graph is acyclic if it has no cycle.

The connected components are the classes of the relation of being joined
by a walk. A graph is connected when it is nonempty and has one component.
A forest is an acyclic graph, and a tree is a connected forest.
The empty graph has no components.
\end{definition}

The walk relation is reflexive by the length-zero walk, symmetric by
reversal, and transitive by concatenation. Any walk can be shortened to a
path with the same endpoints: if \(v_i=v_j\) with \(i<j\), delete the
segment after \(v_i\) through \(v_j\). Every edge after the deletion is
still present, and the length strictly decreases. Thus a nonempty graph
is connected precisely when every pair of vertices is joined by a path.

\begin{theorem}[Trees and spanning forests]\label{p204:W36}
For a nonempty graph on \(n\) vertices, the following conditions are
equivalent:
\begin{enumerate}
\item the graph is a tree;
\item it is connected and has \(n-1\) edges;
\item every pair of vertices has exactly one path between them.
\end{enumerate}
A connected graph has a spanning tree. An arbitrary graph has a
spanning forest that restricts to a spanning tree in each component.
\end{theorem}

\begin{proof}
Two distinct paths with the same endpoints contain a cycle in their
union. To see this, orient both paths from the same endpoint, take the
first vertex where their outgoing edges differ, and follow the second
path from this vertex to its first return to the first path. It cannot
return to the common initial segment because it is a path. This segment
of the second path and the corresponding segment of the first path
form a cycle. Conversely, two arcs of a cycle between distinct vertices
give different paths. Together with the walk reduction, this proves
the equivalence of the first and third conditions.

To construct a spanning tree in a connected graph, start with one
vertex. While the current vertex set is proper, a path from it to an
outside vertex contains an edge leaving it. Add one such edge and its
outside endpoint. The new endpoint has degree one in the growing
subgraph, so this step creates no cycle and preserves connectivity.
There are \(n-1\) steps. A tree has no edge outside the spanning tree
just constructed: any such edge together with the tree path between its
ends forms a cycle. Consequently every tree has \(n-1\) edges, and a
connected graph with \(n-1\) edges consists of its spanning tree.
Apply the construction separately in the components for the last
assertion.
\end{proof}

\begin{theorem}[Leaves and the degree sum]\label{p204:W37}
A tree with at least two vertices has at least two leaves, where a leaf
has degree one. In every finite graph,
\[
  \sum_{v\in V}d_G(v)=2|E|.
\]
In a tree on \(n\) vertices the sum is \(2(n-1)\). In any finite graph
the number of odd-degree vertices is even.
\end{theorem}

\begin{proof}
Choose a path of maximum length in the tree. It has two distinct ends.
A neighbor of an end different from the next path vertex cannot lie
outside the path by maximality, and cannot lie on the path because that
would form a cycle. Each end is therefore a leaf. For the degree sum,
count incident pairs \((v,e)\) first by \(v\) and then by \(e\). Each
edge contributes exactly two. Reduction modulo two gives the parity
claim, and Theorem~\ref{p204:W36} gives the tree specialization.
\end{proof}

\section{Bridges and blocks}

Deleting an edge keeps all vertices. A bridge is an edge whose deletion
increases the number of components. Deleting a vertex also deletes its
incident edges. In a connected graph a cut vertex is a vertex whose
deletion leaves at least two components. A graph is \emph{2-connected}
when it has at least three vertices, is connected, and remains
connected after deletion of any one vertex. A block is a maximal
2-connected subgraph, a bridge with its two endpoints, or a singleton
at an isolated vertex.

\begin{theorem}[Forests and the block--cut tree]\label{p204:W38}
A forest on \(n\) vertices with \(c\) components has \(n-c\) edges.
An edge is a bridge exactly when it lies on no cycle. Every edge belongs
to exactly one block, distinct blocks share at most one vertex, and
a vertex lies in at least two blocks exactly when it is a cut vertex
of its component. For each connected component, the bipartite incidence
graph of its blocks and its cut vertices is a tree. An isolated vertex
gives a single block node.
\end{theorem}

\begin{proof}
Sum the edge counts of the component trees. For the bridge assertion,
if the endpoints of an edge can still be joined after it is deleted,
that path and the deleted edge form a cycle. Replacing each use of the
edge in any old walk by the alternative path shows that no component
splits. If there is no alternative path, the endpoints split.

Two 2-connected subgraphs with two common vertices have 2-connected
union. After one vertex is deleted, each piece remains connected and
they still have a common vertex. Thus distinct maximal 2-connected
subgraphs cannot share two vertices. Each cycle is itself 2-connected
and can be extended to a maximal such subgraph. A 2-connected graph
has no bridge: if a bridge splits it into sides \(A\) and \(B\), then
at least one side has at least two vertices. Deleting the endpoint
of the bridge on that side disconnects the rest of that side from the
other side. These facts give unique membership for every edge, including
the bridge edges. A bridge block cannot share both endpoints with a
2-connected block, since the latter would supply an alternative path
between them. Two different bridge blocks also share at most one
vertex.

For a connected component, temporarily include \emph{all} its vertices
as nodes along with its blocks, joining a vertex node to each block
that contains it. The incidence graph is connected, since each original
edge can be traversed through its block. Suppose it has a cycle with
distinct block nodes \(B_1,\ldots,B_k\) and distinct connecting vertex
nodes. One has \(k\geq 3\), since two blocks cannot share two vertices.
In the union of these blocks, deleting any vertex leaves each piece
connected and nonempty, including the bridge pieces if any. At most
one of the designated links around the cycle disappears. The remaining
links still join all the pieces. The union has at least three vertices
and is 2-connected. This contradicts maximality if there is a
2-connected piece, and contradicts the bridge property if there is a
bridge piece. An isolated-vertex block cannot occur on the cycle.
The incidence graph is therefore a tree.

If a vertex is incident with at least two block nodes, removing its
node disconnects the incidence tree, with other original vertices on
at least two sides. A walk avoiding that vertex in the original graph
would give a walk through blocks avoiding the same node in the tree,
which is impossible between those sides. If the vertex is incident with
a single block, its incidence node is a leaf. For any two remaining
vertices, follow their incidence-tree path and replace each visit
through a block by a walk inside that block with the vertex deleted.
Such walks exist since deleting one vertex in a 2-connected block or
bridge block leaves it connected whenever it has remaining vertices.
If there are fewer than two remaining vertices there is no cut.
Thus exactly the shared vertices are cut vertices. Remove the
non-cut vertex leaves of the incidence tree. The resulting block--cut
graph is a tree, since paths between its retained nodes never used a
deleted leaf.
\end{proof}

\section{Euler trails and Hamilton cycles}

A trail is a walk with distinct edges. It is Eulerian if it uses every
edge. The \emph{edge support} is the subgraph induced by all vertices
of positive degree. A Hamilton cycle is a cycle containing every vertex.

\begin{theorem}[Euler circuits and open trails]\label{p204:W39}
Suppose \(G\) has at least one edge and connected edge support. Then
it has a closed Euler trail if and only if every vertex has even degree.
It has an open Euler trail if and only if it has exactly two vertices
of odd degree, which must be its endpoints. If there are no edges and
at least one vertex, any length-zero walk is a closed Euler trail.
The empty vertex set has no walk under Definition~\ref{p204:W35}.
\end{theorem}

\begin{proof}
In a trail each interior arrival is paired with a departure. Hence the
used degree is even at all vertices for a closed trail and is odd
exactly at the two ends for an open trail. An Euler trail also forces
the edge support to be connected when there are edges.

In the even case, start at a vertex incident with an edge and grow a
trail by unused edges until it cannot be extended at its end. If the
trail ended away from its start, the used degree at the end would be
odd, while its total degree is even. An unused incident edge would
then allow extension. The trail is therefore closed. The unused edges
have even degree at every vertex. If some edge remains, connectivity
of the edge support gives a vertex on the trail incident with an
unused edge. Indeed, if no such vertex existed, all remaining edges
would lie in a different part of the support with no joining edge.
Grow a maximal trail in the unused edges from the shared vertex.
It closes by the same degree argument, and can be inserted into the
original trail at that vertex. Each insertion consumes an edge, so
the process terminates with a closed Euler trail.

In the odd case, start at one odd-degree vertex. A maximal trail
cannot end at a different even-degree vertex by the unused-edge
argument. It cannot close at the starting vertex, since the used degree
there would be even and the total degree odd. It must therefore end at
the other odd-degree vertex. All unused degrees are now even, and the
same insertion of closed trails consumes the remaining edges without
changing the endpoints.
\end{proof}

\begin{theorem}[Ore's and Dirac's conditions]\label{p204:W40}
Let \(n=|V|\geq 3\). If \(d_G(u)+d_G(v)\geq n\) for every pair of
distinct nonadjacent vertices, then \(G\) has a Hamilton cycle.
In particular, this holds if \(d_G(v)\geq n/2\) at every vertex.
\end{theorem}

\begin{proof}
The degree-sum hypothesis forces connectivity. For vertices in distinct
components of sizes \(a\) and \(b\), their degree sum is at most
\(a+b-2\leq n-2\). Choose a longest path \(v_1,\ldots,v_k\).
Since the graph is connected and \(n\geq3\), there is a path on at
least three vertices. All neighbors of either endpoint are on the path.
If \(v_1v_k\) is an edge, the path closes to a cycle on the same vertices.
Otherwise consider the index sets
\[
  I=\{i:1\leq i<k,\ v_1v_{i+1}\in E\},\qquad
  J=\{i:1\leq i<k,\ v_iv_k\in E\}.
\]
They have sizes \(d_G(v_1)\) and \(d_G(v_k)\). Their sizes sum to at least
\(n>k-1\), so they have a common index \(i\). Then
\[
  v_1,\ldots,v_i,v_k,v_{k-1},\ldots,v_{i+1},v_1
\]
is the desired cycle on the path vertices. If this cycle omitted
a vertex, connectivity would give an outside vertex adjacent to
one on the cycle. Starting at that outside vertex and then following
the cycle without closing it gives a longer path. This contradiction
shows that the cycle is Hamiltonian. The minimum-degree hypothesis
implies the required degree sum.
\end{proof}

\section{Counting trees}

\begin{theorem}[Pr\"ufer coding and Cayley's formula]\label{p204:W41}
Let \(V\) be a linearly ordered set of \(n\geq 2\) labels. Trees on \(V\)
are in bijection with sequences of length \(n-2\) over \(V\).
In the corresponding sequence a vertex \(v\) occurs \(d_G(v)-1\) times.
There are \(n^{\,n-2}\) trees on \(V\). On a one-element label set there
is one tree.
\end{theorem}

\begin{proof}
Encode a tree by removing its smallest leaf and recording that leaf's
neighbor, stopping when two vertices remain. Removing a leaf leaves a
tree, since a path between two other vertices cannot use the removed
leaf. For each vertex \(v\), all but one of its incident edges are
removed while \(v\) remains and a leaf neighbor is removed. Each such
removal records \(v\). The last edge is removed when \(v\) is itself
deleted, or is the unrecorded final edge. This proves the occurrence
count. The same argument applies to the residual tree at each stage
and the corresponding suffix of the code. Its leaves are therefore
exactly the residual labels absent from that suffix.

For a sequence \(a_1,\ldots,a_{n-2}\), start with \(V\) as the available
labels. At stage \(i\), let \(b_i\) be the smallest available label
not occurring in \(a_i,\ldots,a_{n-2}\). There are two more available
labels than entries in this suffix, and all suffix entries remain
available by the rule for deletions. Join \(b_i\) to \(a_i\) and remove
\(b_i\) from the available labels. Join the last two labels. In reverse
order this construction starts with one edge and adds one new leaf
at each step, so it produces a tree. On the labels available at stage
\(i\), the constructed residual tree has degree equal to one plus
the number of occurrences in the suffix. This holds for the final
edge. Reversing one step adds a new label of degree one, absent from
the suffix, and increases the degree of the newly recorded label by
one. This proves the formula by induction. It follows that encoding
the constructed tree removes exactly \(b_i\) at stage \(i\) and records
\(a_i\). Conversely, the leaf criterion in the residual tree forces
the decoder to undo each step of any encoding. The bijection gives
the count, including the empty code when \(n=2\).
\end{proof}

\begin{theorem}[Weighted matrix--tree formula]\label{p204:W42}
Let \(G\) have \(n\geq1\) vertices and assign each edge a weight \(w_e\)
in a field \(K\). Define the weighted Laplacian \(L\) by
\[
  L_{uv}=\begin{cases}
    \sum_{e\ni u} w_e & (u=v),\\
    -w_{\{u,v\}} & (u\ne v,\ uv\in E),\\
    0 & (u\ne v,\ uv\notin E).
  \end{cases}
\]
For every vertex \(r\), the determinant of the matrix obtained by
deleting row and column \(r\) equals
\[
   \sum_{T} \prod_{e\in E(T)} w_e ,
\]
summed over the spanning trees. For \(K=\mathbb{Q}\) and all weights one
this is the number of spanning trees. It is zero for a disconnected
graph and is one for a single vertex.
\end{theorem}

\begin{proof}
Orient each edge arbitrarily and form the incidence matrix \(D\), with
\(+1\) at its head and \(-1\) at its tail in each column. Then
\(L=D\operatorname{diag}(w_e)D^{\mathsf T}\), by reading a column's two
nonzero entries. Let \(D_r\) omit row \(r\) and write \(D_{r,S}\) for
its columns in a subset \(S\) of the edges, with the inherited ordering.
Expanding by multilinearity in the columns yields
\[
 \det\bigl(D_r\operatorname{diag}(w_e)D_r^{\mathsf T}\bigr)
   =\sum_{|S|=n-1}\det(D_{r,S})^2\prod_{e\in S} w_e .
\]
For clarity, in the determinant of a product \(AB\) of sizes \(p\times m\)
and \(m\times p\), expand each output column as the corresponding linear
combination of columns of \(A\). A repeated input column gives zero.
For each \(p\)-element subset, group the terms over its permutations.
The determinant sign is the permutation sign, and the sum of products
of coefficients is exactly the determinant of the corresponding rows
of \(B\). Taking \(A=D_r\) and \(B=\operatorname{diag}(w_e)D_r^{\mathsf T}\)
gives the displayed expansion.

If \(S\) is a spanning tree and \(n>1\), choose a leaf different from
\(r\) by Theorem~\ref{p204:W37}. Its row in \(D_{r,S}\) has one nonzero
entry \(\pm1\). Expanding in that row reduces to the same determinant
for the tree with that leaf and edge removed. The empty determinant
is one, so the determinant is \(\pm1\). If \(S\) is not a spanning
tree, the graph \((V,S)\) is disconnected by Theorem~\ref{p204:W36}.
Sum the rows in a component not containing \(r\). The sum is zero
since every column with an endpoint in the component has both
endpoints there. The row dependence is nontrivial. Thus only spanning
trees contribute to the expansion. The sole empty tree on a single
vertex has empty product one, and a disconnected graph has no spanning
tree.
\end{proof}

\section{Fundamental cycles and cuts}

Identify a set of edges with its indicator in the vector space
\(\mathbb{F}_2^E\), with addition given by symmetric difference. For a
vertex set \(S\), let \(\delta(S)\) be the edges with exactly one
endpoint in \(S\). The cut space is the span of these vectors. The cycle
space is the span of the edge sets of simple cycles.

\begin{theorem}[Bases of the cycle and cut spaces]\label{p204:W43}
Let \(G\) have \(n\) vertices, \(m\) edges, and \(c\) components.
Fix a spanning forest \(F\) with one tree in each component.
For an edge \(e\notin F\), adjoining \(e\) to its forest path gives
its fundamental cycle. For \(e\in F\), deleting \(e\) splits its
forest tree and one side determines its fundamental cut. These
fundamental cycles and cuts are bases of the cycle and cut spaces
respectively. The dimensions are \(m-n+c\) and \(n-c\), and the two
spaces are orthogonal under the coordinate dot product over
\(\mathbb{F}_2\). The cycle space consists of exactly the even-degree
edge sets.
\end{theorem}

\begin{proof}
The forest path exists because \(F\) spans each component. The
fundamental cycle has only one nonforest edge, namely \(e\). Thus
the cycles are independent by their nonforest coordinates. Every
sum of cycles gives even degree at each vertex. Conversely, start
with an even-degree edge set, and add the fundamental cycle for each
nonforest edge present. The remainder is an even-degree subset of
the forest, necessarily empty: a nonempty edge subset of a forest
has a component with an edge and hence a leaf. Thus the fundamental
cycles span all even-degree sets.

The fundamental cut of a forest edge \(e\) has only that edge among
its forest coordinates, because the two sides are components of
the forest tree after deleting \(e\). Hence the fundamental cuts
are independent. The identity
\(\delta(S)+\delta(T)=\delta(S\mathbin{\triangle}T)\) follows by
comparing the endpoint indicators modulo two. Thus every cut-space
vector is itself \(\delta(S)\) for some \(S\). Adding the fundamental
cuts at its nonzero forest coordinates leaves a cut with no forest
edge. Membership in its defining vertex set is constant along
each forest tree, hence constant on each component of \(G\).
That cut is empty, so the fundamental cuts span. By
Theorem~\ref{p204:W38} there are \(n-c\) forest edges and \(m-n+c\)
nonforest edges. Finally, for an even-degree edge set \(H\),
the number of its edges in \(\delta(S)\) is congruent modulo two to
\(\sum_{v\in S}d_H(v)\): internal edges contribute twice. This is
zero, proving orthogonality.
\end{proof}

The finite graph conventions and classical statements in this chapter
may be compared with \cite[Chapters 3--5, 18, and 20]{BM}.
The determinant expansion uses the same Cauchy--Binet interface as
\cite{D05}.
\chapter{Network flows, integrality, and Menger's theorem}
\label{p204:ch07}

\section{Finite networks and the cut inequality}

\begin{definition}[Capacitated directed network]\label{p204:W44}
A finite network consists of a finite vertex set \(V\), a finite arc
set \(A\), tail and head maps \(\tau,\eta:A\to V\), finite capacities
\(c_a\in[0,\infty)\), and distinct vertices \(s,t\in V\).
Parallel and oppositely directed arcs are allowed. They retain their
arc identities. A flow assigns numbers \(f_a\) with \(0\leq f_a\leq c_a\)
and with divergence
\[
  \operatorname{div}_f(v)
     = \sum_{\tau(a)=v} f_a-\sum_{\eta(a)=v}f_a
     =0 \qquad (v\notin\{s,t\}).
\]
The value is \(|f|=\operatorname{div}_f(s)\). Here \(|f|\) denotes the
signed value, rather than its absolute value. An \(s\)--\(t\) cut is a
set \(S\subseteq V\) containing \(s\) and excluding \(t\). Write
\[
 A^+(S)=\{a:\tau(a)\in S,\ \eta(a)\notin S\},\qquad
 A^-(S)=\{a:\tau(a)\notin S,\ \eta(a)\in S\}.
\]
Its capacity is \(c(S)=\sum_{a\in A^+(S)}c_a\).
\end{definition}

Every arc contributes zero to the sum of the divergences over \(V\),
so \(\operatorname{div}_f(t)=-|f|\). A loop makes no contribution to any
divergence or cut. One can remove loops for the optimization problem
and extend flows back by zero on the removed arcs. We allow loops in
circulation decompositions.

\begin{proposition}[Weak flow--cut duality]\label{p204:W45}
For every feasible flow \(f\) and \(s\)--\(t\) cut \(S\),
\[
 |f|=\sum_{a\in A^+(S)}f_a-\sum_{a\in A^-(S)}f_a\ \leq\ c(S).
\]
Equality holds exactly when every outgoing cut arc is saturated
and every incoming cut arc has zero flow.
\end{proposition}

\begin{proof}
Sum the divergence over vertices in \(S\). An arc with both endpoints
in \(S\) contributes once with each sign, so it cancels. The only
nonzero divergence in \(S\) can occur at \(s\). The difference
\(c(S)-|f|\) is now a sum of outgoing slacks \(c_a-f_a\) and incoming
flows \(f_a\), all nonnegative. It is zero precisely when each term
is zero.
\end{proof}

\section{Residual paths and real capacities}

For a feasible \(f\), give each original arc \(a:u\to v\) two possible
residual arcs. Its forward copy is \(u\to v\) with residual capacity
\(c_a-f_a\), and its backward copy is \(v\to u\) with residual capacity
\(f_a\). Use only positive residual capacities when defining the
residual digraph. The copies are tagged by \(a\) and by the direction
of modification, so even when endpoints coincide they remain separate.

\begin{lemma}[Augmentation and the reachable cut]
\label{p204:residual-path}
If there is a directed residual walk from \(s\) to \(t\), the value
can be increased. Otherwise the vertices reachable from \(s\) in the
residual digraph give a cut of capacity equal to the flow value.
\end{lemma}

\begin{proof}
Shorten the walk to one with distinct vertices by removing segments
between repeated vertices, retaining the direction of every remaining
arc. Let \(\varepsilon>0\) be the minimum residual capacity on this
path. Increase \(f_a\) by \(\varepsilon\) for each forward copy used
and decrease it by \(\varepsilon\) for each backward copy used. A
path with distinct vertices uses at most one copy of any original
arc: use of both would revisit an endpoint, and a loop cannot occur.
The bounds on every original arc are therefore preserved by the
choice of \(\varepsilon\). Each residual arc increases divergence by
\(\varepsilon\) at its tail and decreases it by \(\varepsilon\) at its
head, even when it is a backward copy. The changes cancel at every
interior vertex of the path and increase the value by \(\varepsilon\).

If no such walk exists, the reachable set \(S\) contains \(s\) and
excludes \(t\). A positive forward residual copy of an arc in \(A^+(S)\)
would leave \(S\), so every such original arc is saturated. A positive
backward copy of an arc in \(A^-(S)\) would also leave \(S\), so every
such original arc has zero flow. Proposition~\ref{p204:W45} gives the
claimed equality.
\end{proof}

\begin{theorem}[Maximum flow and minimum cut for real capacities]
\label{p204:W46}
Every finite network has a maximum-value real flow and a minimum
\(s\)--\(t\) cut. Their values are equal.
\end{theorem}

\begin{proof}
The feasible set is nonempty because it contains zero. Every flow
value is bounded above by the sum of capacities of the arcs out of
\(s\). Choose feasible flows whose values converge to the supremum.
For each arc in turn, take a subsequence on which the flow coordinate
converges in its closed bounded interval \([0,c_a]\). The needed
one-dimensional fact follows by repeatedly bisecting a bounded
interval, retaining a closed half with infinitely many terms, and
choosing successively later terms in the nested retained halves.
Their lengths tend to zero, giving a convergent subsequence by
completeness of the reals. Since there are finitely many arcs, the
last subsequence converges on all coordinates. The limit remains
feasible by the closed bounds and the finite linear conservation
equalities. Its value is the supremum by the finite value sum.
The case of no arcs uses just the zero flow.

The residual graph of a maximizing flow has no \(s\)-to-\(t\)
walk by Lemma~\ref{p204:residual-path}. The reachable cut has capacity
equal to its value, and Proposition~\ref{p204:W45} shows that this cut
is minimum. This argument uses attainment in the finite-dimensional
feasible set, without assuming termination of arbitrary real
augmentations.
\end{proof}

\begin{theorem}[Residual optimality and path--cycle decomposition]
\label{p204:W47}
A feasible flow is maximum exactly when it has no \(s\)-to-\(t\)
residual path. A nonnegative circulation on a finite directed
multigraph is a nonnegative sum of directed simple cycles, allowing
a loop as a one-arc cycle. A feasible flow of positive value is
a sum of such cycles and simple directed \(s\)-to-\(t\) paths, with
path weights summing to its value. At value zero only cycles are
needed, and negative value gives paths in the opposite direction
of total weight equal to the opposite of the value. Each decomposition
preserves every original arc load. If the loads are integers, the
summands can be taken with positive integer weights.
\end{theorem}

\begin{proof}
A residual path strictly improves the value by
Lemma~\ref{p204:residual-path}. In its absence the cut from the lemma
and the weak dual bound certify optimality.

A circulation is a nonnegative assignment with zero divergence at all
vertices. If its positive support has any arc, follow positive arcs.
At each head reached there is a positive outgoing arc by conservation.
Since there are finitely many vertices, the first repetition gives
a directed cycle with distinct vertices apart from its closure.
Subtract the minimum load on its arcs along the cycle. This preserves
nonnegativity and conservation, and removes at least one arc from the
positive support. Iteration stops after finitely many subtractions.
Adding them back gives the cycle decomposition. All subtractions
are integral when the initial loads are.

For a flow of positive value \(v\), add a new distinguished arc
\(t\to s\) of load \(v\). The enlarged assignment is a circulation.
Every cycle using the new arc gives, on deleting that arc, a simple
directed \(s\)-to-\(t\) path. Other cycles are kept as they were.
The sum of path weights is the total load \(v\) on the new arc.
The original arc loads are unchanged in the resulting decomposition.
For value zero the flow is already a circulation. For negative value
interchange \(s\) and \(t\) and apply the positive case.
\end{proof}

\section{Integral flows and disjoint paths}

\begin{theorem}[Integrality and both finite Menger equalities]
\label{p204:W48}
If all capacities are nonnegative integers, some maximum flow has
integer arc loads. For distinct vertices \(s,t\) of a finite simple
undirected graph \(G\), the maximum number of pairwise edge-disjoint
\(s\)-to-\(t\) paths equals the minimum size of an edge set whose
deletion separates \(s\) and \(t\). If \(s,t\) are also nonadjacent,
the maximum number of pairwise internally vertex-disjoint
\(s\)-to-\(t\) paths equals the minimum size of a set of vertices
in \(V\setminus\{s,t\}\) whose deletion separates them.
\end{theorem}

\begin{proof}
Begin with the zero flow. At each step use any residual path and
augment by its minimum positive residual capacity. The flow remains
integral, and its value grows by at least one. Since the sum of
capacities out of \(s\) is a finite upper bound, after finitely many
steps no residual path remains. Theorem~\ref{p204:W47} then gives
an integral maximum.

For the edge statement, replace each undirected edge by two opposite
arcs of capacity one. For an \(s\)--\(t\) cut \(S\), its capacity is
exactly the size of \(\delta(S)\) in the undirected graph. Such
an edge cut separates \(s,t\). Conversely, if an edge set separates
them, take \(S\) to be the vertices reachable from \(s\) after that
deletion. Then \(\delta(S)\) is contained in the deleted set.
Thus the minimum cut capacity equals the minimum separator size.

Undirected edge-disjoint paths, oriented from \(s\) to \(t\), give
an integral flow of their number. Conversely, take an integral
maximum and, for each undirected edge, subtract the smaller of its
two opposite loads from both. This preserves divergence and all
capacity bounds and leaves at most one positive load on that edge.
Decompose the result as in Theorem~\ref{p204:W47}. Every path uses at
least one arc and all arc loads are at most one, so every path has
weight one. No two paths can use the same undirected edge, since
it now carries load at most one in total. At value zero the empty
path family suffices. The real flow--cut equality therefore gives
the edge statement with integer path counts.

For the vertex statement, write \(n=|V|\). Split every internal vertex
\(v\) into \(v^{\mathrm{in}},v^{\mathrm{out}}\), with an arc
\(v^{\mathrm{in}}\to v^{\mathrm{out}}\) of capacity one. Do not split
\(s,t\). For every undirected edge \(uv\) put a transit arc from
\(u^{\mathrm{out}}\) to \(v^{\mathrm{in}}\) and one from
\(v^{\mathrm{out}}\) to \(u^{\mathrm{in}}\), each of capacity \(n\),
using the unsplit vertex at an endpoint. Every simple directed
\(s\)-to-\(t\) path in this network projects to a simple undirected
path: at an internal vertex the only way to pass from its input to
its output is its split arc, and a repetition in the projection would
repeat that arc's input and output. Conversely an undirected
\(s\)-to-\(t\) path lifts to a directed one.

Deleting the split arcs of a vertex separator prevents all directed
\(s\)-to-\(t\) paths, so the resulting reachable cut in the original
network has outgoing arcs among those split arcs. In particular,
deleting all \(n-2\) split arcs prevents such a path, since \(s,t\)
are nonadjacent. The minimum cut has capacity at most \(n-2\)
and hence contains no outgoing transit arc of capacity \(n\).
Its outgoing split arcs give a vertex separator, since every lifted
path must cross the cut. This proves equality of the minimum cut
capacity and the minimum vertex separator size.

Internally disjoint undirected paths give a feasible lifted integral
flow: the unit split arcs are not shared, and there are at most
\(n-2\) paths because each contains an internal vertex. Conversely
an integral maximum decomposes into paths and cycles. Every lifted
path passes through at least one split arc, forcing its weight to
be one. No two such paths share a split arc, so their projections
have disjoint interiors. This proves the vertex equality.
\end{proof}

\begin{corollary}[The three connectivity bounds]
\label{p204:whitney}
Let \(G\) be a finite simple graph on \(n\geq2\) vertices. Let
\(\kappa\) be the least number of vertices whose removal either
disconnects the remaining graph or leaves one vertex, and let
\(\lambda\) be the least number of edges whose removal disconnects
the graph. Set both to zero for an initially disconnected graph.
If \(\delta\) is the minimum degree, then
\[
  \kappa\leq\lambda\leq\delta.
\]
For a complete graph, \(\kappa=n-1\).
\end{corollary}

\begin{proof}
Removing all edges at a minimum-degree vertex gives the upper
bound on \(\lambda\). In a connected graph take a minimum edge
cut and a partition \(S,V\setminus S\) induced by a component
after the cut is removed. Both sides are nonempty and every edge
between them was removed. By minimality these crossing edges
form a cut of size \(\lambda\). If every cross-pair is adjacent,
its size is \(|S|(n-|S|)\geq n-1\), since the two positive side
sizes \(a,b\) satisfy \(ab-a-b+1=(a-1)(b-1)\geq0\). The bound
\(\kappa\leq n-1\) follows by retaining just one vertex. Otherwise
choose nonadjacent \(u,v\) on opposite sides. Choose one endpoint
different from \(u,v\) of each crossing edge, possible because
\(uv\) is absent. Deleting the chosen vertices removes all crossing
edges and leaves both \(u\) and \(v\), so it disconnects the graph.
At most \(\lambda\) vertices have been chosen. The conventions give
the disconnected case and the asserted complete-graph value.
\end{proof}

For historical and textbook comparison, see \cite[Chapters 7 and 9]{BM}
and \cite[Chapters 10 and 15]{Schrijver}.
\chapter{Matchings, covers, and chain decompositions}
\label{p204:ch08}

\section{Matchings and augmenting paths}

\begin{definition}[Matchings and vertex covers]\label{p204:W49}
In a finite simple graph, a matching is a set of edges no two of
which share an endpoint. A vertex is saturated if it is incident
with a matching edge. A matching is perfect if it saturates all
vertices. A vertex cover is a set meeting each edge. In particular,
the empty matching is perfect on the empty graph.
\end{definition}

Every matching has size at most that of every vertex cover.
For each matching edge choose a covered endpoint. Distinct
matching edges give distinct chosen endpoints, so this is an
injection into the cover.

An augmenting path for a matching \(M\) has distinct vertices,
starts and ends at unsaturated vertices, and alternates between
edges outside and inside \(M\), starting and ending outside \(M\).

\begin{theorem}[Berge's augmenting-path criterion]\label{p204:W50}
A matching \(M\) in a finite graph has maximum cardinality if and
only if it has no augmenting path.
\end{theorem}

\begin{proof}
On an augmenting path, remove its \(M\)-edges and insert its other
edges. At each interior vertex one matching edge is replaced by
one, and at the two ends a new vertex is saturated. The result is
a matching with one more edge.

Conversely, suppose a matching \(N\) is larger than \(M\). In the
graph with edges in \(M\mathbin{\triangle}N\), each vertex has
degree at most two, and at a degree-two vertex one edge belongs
to each matching. A finite connected graph of maximum degree at
most two with edges is a path or a cycle: if it has a degree-one
vertex, follow its edges starting there without retracing. The
walk cannot enter a cycle, since the entry would require three
incident edges or would return to the degree-one start. It
ends at another degree-one vertex and exhausts the component.
If every degree is two, following the edges instead closes a cycle
and exhausts the component. Hence the nontrivial components of
the symmetric difference are alternating paths and alternating
even cycles. Some component has more \(N\)-edges than \(M\)-edges
because the total difference of sizes is positive. It must be
a path that begins and ends with an \(N\)-edge. Its endpoints
are unsaturated in \(M\): an incident edge of \(M\setminus N\)
would extend this component, and an incident common edge would
give two \(N\)-edges at the endpoint. This path augments \(M\).
\end{proof}

\section{Bipartite matching and covering}

In a bipartite graph let the vertex parts be the disjoint sets
\(L\) and \(R\), and let \(N(S)\subseteq R\) denote the neighbors
of a subset \(S\subseteq L\).

\begin{lemma}[Alternating reachability from unmatched left vertices]
\label{p204:alternating-reach}
For a maximum matching \(M\), let \(U\subseteq L\) be the unmatched
left vertices. Direct every nonmatching edge from left to right,
and every matching edge from right to left. Let \(S\subseteq L\)
and \(T\subseteq R\) be the vertices reachable in this digraph
from \(U\), including \(U\) itself. Then
\[
 T=N(S),\qquad \text{matching partners give a bijection }T\simeq S\setminus U .
\]
\end{lemma}

\begin{proof}
A reached unmatched right vertex would give a directed path from
\(U\) to it. Shorten a directed walk by removing repeated-vertex
segments if necessary. Since the only directed right-to-left edges
are matching edges, the path is augmenting, contrary to
Theorem~\ref{p204:W50}. Thus every \(y\in T\) is matched, and its
partner is in \(S\setminus U\) by the matching edge directed out
of \(y\). Conversely a matched \(x\in S\) is not a root of the
search. Any directed path to it has positive length and must end
with the arc from its matching partner, so that partner lies in
\(T\). This proves the bijection. Every \(y\in T\) was reached by
an edge from some \(x\in S\), giving \(T\subseteq N(S)\).
For \(x\in S\), each nonmatching edge leads directly to its right
endpoint, and its matching edge, if present, leads in the undirected
graph to the partner just proved to lie in \(T\). Thus \(N(S)\subseteq T\).
\end{proof}

\begin{theorem}[Hall's marriage theorem]\label{p204:W51}
A finite bipartite graph has a matching saturating \(L\) if and only
if \(|N(X)|\geq |X|\) for every \(X\subseteq L\).
\end{theorem}

\begin{proof}
A matching saturating \(L\) sends distinct vertices of \(X\) to
distinct vertices of \(N(X)\). For sufficiency choose a maximum
matching, which exists because the graph is finite. If its unmatched
left set \(U\) is nonempty, Lemma~\ref{p204:alternating-reach} gives
\[
 |N(S)|=|T|=|S|-|U|<|S|,
\]
contradicting the hypothesis. Thus the matching saturates \(L\).
\end{proof}

\begin{theorem}[K\H{o}nig's matching--cover equality]\label{p204:W52}
In a finite bipartite graph the largest matching size equals the
smallest vertex cover size.
\end{theorem}

\begin{proof}
For a maximum matching use the sets \(S,T,U\) from
Lemma~\ref{p204:alternating-reach}. The set
\((L\setminus S)\cup T\) covers all edges: an edge whose left
endpoint lies in \(S\) has right endpoint in \(N(S)=T\).
Its size is
\[
 |L|-|S|+|T|=|L|-|U|=|M|.
\]
The matching-to-cover injection following Definition~\ref{p204:W49}
gives the reverse bound. This argument includes empty parts.
\end{proof}

\section{Perfect and maximum matchings in general graphs}

For \(S\subseteq V\), write \(G-S\) for the induced graph on
\(V\setminus S\), and \(o(G-S)\) for the number of its odd-order
components.

\begin{theorem}[Tutte's perfect-matching criterion]\label{p204:W53}
A finite graph has a perfect matching if and only if
\(o(G-S)\leq |S|\) for every \(S\subseteq V\).
\end{theorem}

\begin{proof}
In a perfect matching each odd component of \(G-S\) has some
vertex matched outside that component. Such an edge meets \(S\),
and the vertices of \(S\) used by different components are distinct.
This proves necessity.

For sufficiency suppose the inequalities hold and there is no perfect
matching. Adding edges cannot increase the number of odd components
after any prescribed deletion: each new edge either has no effect
or merges two components, and the parity of the union is the sum
of their parities modulo two. We can therefore enlarge the graph
on its fixed finite vertex set until it is edge-maximal without a
perfect matching. Every missing edge added alone now creates a
perfect matching, and the inequalities still hold.

Let \(U\) be the vertices adjacent to all other vertices of this
enlarged graph. We claim that every component of \(G-U\) is complete.
If one is not, take a shortest path between two nonadjacent
vertices in it. Its first three vertices \(x,y,z\) satisfy
\(xy,yz\in E\) and \(xz\notin E\). The middle vertex \(y\) lies
outside \(U\), so it has a distinct nonneighbor \(w\). This \(w\)
is different from \(x,z\). By maximality there is a perfect matching
\(M_1\) of \(G+xz\) containing \(xz\), and a perfect matching \(M_2\)
of \(G+yw\) containing \(yw\). Their symmetric difference consists
of alternating cycles, since all vertices are saturated in both
matchings. If the two added edges are on different cycles,
switch from \(M_1\) to \(M_2\) on the cycle containing \(xz\).
This removes \(xz\) without using \(yw\), giving a perfect matching
of \(G\), a contradiction.

If both added edges lie on one cycle, remove them from that cycle.
The two resulting paths each join an endpoint of \(xz\) to an
endpoint of \(yw\). Each has even length, since the first edge
after the removed \(M_1\)-edge belongs to \(M_2\), and the last
edge before the removed \(M_2\)-edge belongs to \(M_1\).
If \(y\) is paired by these paths with \(x\), use the real edge
\(yz\) to match \(y\) and \(z\). If \(y\) is paired with \(z\),
use \(yx\). In either case one endpoint is thereby removed from
each even-length path. Each remaining path has an even number
of vertices, which can be matched in consecutive pairs along its
edges. Together with the chosen real edge this matches the entire
cycle in \(G\). Use \(M_1\) off that cycle, where all of its edges
are real. Again there is a perfect matching of \(G\). These
contradictions prove the claim.

By the inequality for \(S=U\), assign to each odd component of
\(G-U\) a different vertex of \(U\). Match it to one vertex of
that component, using its universal adjacency. The remaining
vertices in every component of \(G-U\) have even size and can
be paired within their complete component. The inequality for
\(S=\varnothing\) ensures that all components of \(G\) have even
size and hence \(|V|\) is even. The unused vertices of \(U\)
therefore also have even size, and can be paired in the complete
subgraph on \(U\). This final perfect matching is a contradiction.
The original graph must have had a perfect matching.
\end{proof}

\begin{theorem}[The Tutte--Berge formula]\label{p204:W54}
The least number of unmatched vertices in a matching of a finite
graph \(G\) is
\[
 d=\max_{S\subseteq V}\bigl(o(G-S)-|S|\bigr).
\]
Equivalently the maximum matching size is \((|V|-d)/2\).
\end{theorem}

\begin{proof}
For a matching and a fixed \(S\), at most \(|S|\) odd components
of \(G-S\) have a matching edge to \(S\). Each other odd component
has an unmatched vertex. Thus at least \(d\) vertices must remain
unmatched. The maximum defining \(d\) exists and \(d\geq0\) by the
choice \(S=\varnothing\). Also \(o(G-S)\equiv |V|-|S|\pmod 2\),
because even components contribute zero modulo two to the order.
Every term in the maximum, and hence \(d\), has parity \(|V|\).

Adjoin a set \(D\) of \(d\) new vertices, each adjacent to every
other vertex, to form \(H\). For a deletion set \(X\subseteq V(H)\)
that contains \(D\), the definition of \(d\) gives
\[
 o(H-X)=o\bigl(G-(X\cap V)\bigr)\leq |X\cap V|+d=|X|.
\]
If some new vertex remains, the remaining graph is connected and
has at most one odd component. This gives the inequality for
nonempty \(X\). For \(X=\varnothing\) its order \(|V|+d\) is even,
so the inequality also holds. By Theorem~\ref{p204:W53}, \(H\)
has a perfect matching. Delete its new vertices and the matching
edges incident with them. At most \(d\) original vertices lose
their partners. The lower bound already proved is therefore
attained, and the size formula follows.
\end{proof}

\section{Chains and antichains}

A chain in a poset is a subset whose elements are pairwise comparable,
and an antichain is a subset with no two distinct comparable elements.
The width is the maximum size of an antichain and the height is the
maximum size of a chain. Both are zero for the empty poset.
Decompositions below are partitions into nonempty parts.

\begin{theorem}[Dilworth's and Mirsky's equalities]\label{p204:W55}
For a finite poset \(P\), the minimum number of chains in a chain
decomposition equals the width. The minimum number of antichains
in an antichain decomposition equals the height.
\end{theorem}

\begin{proof}
Put \(n=|P|\) and form the bipartite graph with left and right
copies of \(P\), joining \(x_L\) to \(y_R\) if \(x<y\).
Let \(m\) be its maximum matching size. Read a matching as
arrows \(x\to y\) on \(P\). Each vertex has at most one incoming
and one outgoing arrow. Every arrow is strictly increasing,
so there is no directed cycle. The arrows therefore give
vertex-disjoint directed paths, including isolated vertices.
Indeed, following predecessors must stop at a vertex with no
predecessor, after which following successors exhausts the path.
Transitivity makes the vertices on each path a chain. A matching
of size \(m\) thus gives \(n-m\) chains, since a path on \(k\)
vertices has \(k-1\) arrows. Conversely a partition into \(b\)
chains supplies a matching of size \(n-b\) by joining successive
vertices in each chain from the left to the right copy. This
proves that the minimum chain count is \(n-m\).

By Theorem~\ref{p204:W52} this bipartite graph has a vertex cover
of size \(m\). Take all elements of \(P\) neither of whose copies
is in the cover. At most \(m\) elements are excluded, and no
two retained elements satisfy \(x<y\), since that would leave
\(x_Ly_R\) uncovered. Thus there is an antichain of size at least
\(n-m\). Every antichain meets each part of any chain decomposition
in at most one element, proving the opposite bound on width.

For the second assertion let \(h\) be the height. A chain of
size \(h\) meets each antichain in at most one element, giving
a lower bound of \(h\) on the antichain count. For each \(x\),
define \(\ell(x)\) to be the largest size of a chain with largest
element \(x\). If \(x<y\), a longest such chain at \(x\) can be
extended by \(y\), so \(\ell(y)>\ell(x)\). Hence each level of
\(\ell\) is an antichain. The nonempty levels partition \(P\)
into at most \(h\) antichains. The empty poset uses the empty
partition in both statements.
\end{proof}

For the classical matching and decomposition statements, see
\cite[Chapter 16]{BM} and \cite[Sections 14.2, 16.2, and 24.1]{Schrijver}.
\chapter{Coloring and plane graphs}\label{p204:ch09}

Throughout this chapter \(G=(V,E)\) is a finite simple undirected
graph unless parallel edges are explicitly allowed. We write
\([k]=\{1,\ldots,k\}\) and \([0]=\varnothing\).

\section{Vertex colorings and the degree bounds}

\begin{definition}[Proper coloring and chromatic number]\label{p204:W56}
A proper \(k\)-coloring is a function \(c:V\to[k]\) satisfying
\(c(u)\ne c(v)\) for every \(uv\in E\). The chromatic number
\(\chi(G)\) is the least \(k\) for which such a function exists.
\end{definition}

Assigning a different color to each vertex gives a coloring, so the
minimum exists. The empty graph has chromatic number zero. A nonempty
edgeless graph has chromatic number one, since a constant color works
and no function from a nonempty set to \([0]\) exists. Relabelling
vertices by a graph isomorphism transports a coloring by composition
with its inverse, and relabelling colors by a bijection preserves
inequality of colors. Thus the definition depends only on the graph
and the cardinality of the color set.

\begin{theorem}[Greedy coloring and degeneracy]\label{p204:W57}
A graph of maximum degree \(\Delta\) satisfies \(\chi(G)\le\Delta+1\),
where \(\Delta=0\) for the empty graph. If every nonempty induced
subgraph has a vertex of degree at most \(d\), then \(\chi(G)\le d+1\).
The latter condition is equivalent to the existence of an order in
which each vertex has at most \(d\) neighbors later in the order.
\end{theorem}
\begin{proof}
For the first bound color the vertices in any order. At the time a
vertex is considered, at most \(\Delta\) of its neighbors have already
received colors. Some color among \(\Delta+1\) is therefore available,
and assigning it preserves the condition on all edges already colored
at both ends. For the second bound repeatedly remove a vertex of degree
at most \(d\) in the current induced subgraph. Its later neighbors
in the removal order are precisely its neighbors still present.
Conversely, for any nonempty induced subgraph, the first vertex of
that subgraph in the proposed order has at most \(d\) neighbors in
the subgraph. Coloring in the reverse removal order gives the
claimed bound. All assertions hold for the empty order.
\end{proof}

\begin{theorem}[Bipartiteness and odd cycles]\label{p204:W58}
A finite simple graph is bipartite if and only if it contains no
cycle of odd length. In particular a nonempty bipartite graph admits
a proper two-coloring.
\end{theorem}
\begin{proof}
In a bipartite graph every step of a walk switches parts, so a closed
walk has even length. Conversely, in each connected component choose
a root and color each vertex by the parity of its shortest-path
distance to that root. If the endpoints of an edge had the same
parity, the two root paths together with the edge would form a
closed walk of odd length. Every odd closed walk contains an odd
cycle. To see this, take a shortest odd closed walk in the graph
traced by the original walk. If a vertex occurs twice apart from
the required first and last occurrences, split the walk at that
repetition into two shorter closed walks. Their lengths add to an
odd integer, so one is odd, a contradiction. The shortest walk
therefore has distinct internal vertices. Its length is at least
three because the graph has no loops, and it is a cycle.
Thus the distance parities give a proper bipartition.
\end{proof}

\begin{lemma}[Two initial vertices for a greedy order]
\label{p204:brooks-order}
Suppose \(G\) is a two-connected, noncomplete graph of minimum degree
at least three. Some vertex \(v\) has two nonadjacent neighbors
\(a,b\) such that \(G-\{a,b\}\) is connected.
\end{lemma}
\begin{proof}
If deleting any two vertices leaves a connected graph, take a shortest
path between two nonadjacent vertices and use its first three
vertices. The endpoints of this three-vertex segment are nonadjacent:
a chord would shorten the original path. Its middle vertex is \(v\).

Otherwise choose a pair \(v,w\) such that \(H=G-v\) is connected
and \(H-w\) is disconnected. The block-cut tree of \(H\) from
Theorem~\ref{p204:W38} has at least two leaf blocks. Each meets
the rest of \(H\) in exactly one cut vertex \(t\). Its interior is
the nonempty set of its vertices other than \(t\). No edge of \(H\)
joins that interior to a vertex outside the block: each interior
vertex belongs to just that block, and every edge lies in a block.
There must be a neighbor of \(v\) in each such interior. Otherwise
deleting \(t\) from \(G\) would separate this interior from \(v\),
contrary to two-connectivity. Choose \(a,b\) in the interiors of
two different leaf blocks. They are nonadjacent by the same
edge-in-a-block property. Deleting a vertex from the interior of a
two-connected leaf block leaves a connected graph containing its
cut vertex. For a leaf block consisting of a single bridge the
cut vertex itself remains. Thus \(H-\{a,b\}\) stays connected:
one may retain the whole block-cut tree after removing the two
leaf interiors, and the remaining parts of those blocks still
attach to their cut vertices. Since \(v\) has degree at least
three it has a neighbor other than \(a,b\). Adding it back
preserves connectivity.
\end{proof}

\begin{theorem}[Brooks' theorem]\label{p204:W59}
If a finite connected simple graph is neither a complete graph nor
an odd cycle, then \(\chi(G)\le\Delta(G)\).
\end{theorem}
\begin{proof}
We use the following ordering observation. In a connected graph of
maximum degree at most \(\Delta\), if a specified vertex has degree
less than \(\Delta\), root a spanning tree at that vertex and color
by decreasing tree distance. Every nonroot vertex has an uncolored
parent at its turn, and the root sees fewer than \(\Delta\) neighbors.
Consequently \(\Delta\) colors suffice.

For \(\Delta\le2\), a connected graph is a single vertex, a path,
or a cycle. Indeed a maximal simple path cannot miss an attached
vertex without either extending at an endpoint or giving an
internal vertex degree at least three. If the graph contains a
cycle, its vertices already have degree two and connectivity
excludes any other vertices. Under the stated exclusions the only
possibilities are paths of maximum degree two and even cycles;
alternating colors suffice. Now suppose \(\Delta\ge3\).
If \(G\) is not regular, apply the ordering observation at a vertex
of smaller degree. If \(G\) has a cut vertex \(v\), let \(H_i\)
be the induced graph on \(v\) together with one component of \(G-v\).
Each \(H_i\) is connected, and \(v\) has a neighbor in every such
component. Its degree in each \(H_i\) is strictly less than \(\Delta\).
Color each \(H_i\) using the observation, then permute its color
labels so the colors at \(v\) agree. These colorings combine
because there is no edge between distinct components of \(G-v\).

The remaining graph is two-connected and \(\Delta\)-regular.
Use Lemma~\ref{p204:brooks-order} to choose \(a,b,v\).
Color \(a\) and \(b\) first with the same color. Root a spanning
tree of \(G-\{a,b\}\) at \(v\) and color its other vertices by
decreasing distance to \(v\). Each sees an uncolored parent at
its turn and therefore has an available color. At \(v\), the two
neighbors already assigned the same color guarantee that its
\(\Delta\) neighbors use at most \(\Delta-1\) colors.
This completes the coloring.
\end{proof}

These degree bounds concern the largest local degree or a removal
order. They apply without any choice of a drawing of the graph
\cite[Chapter 14]{BM}\cite[Chapter 5]{Diestel}.

\section{The chromatic polynomial}

For a set of edges \(A\subseteq E\), write \(c(A)\) for the number
of connected components of \((V,A)\), including isolated vertices.
For a simple edge \(e=uv\) we write \(G/e\) for the simple graph
obtained by identifying \(u,v\), removing \(e\), and merging any
parallel edges created by the identification.

\begin{theorem}[Chromatic polynomial and deletion--contraction]
\label{p204:W60}
There is a unique polynomial \(P_G(q)\in\mathbb Z[q]\) whose value
at each nonnegative integer \(k\) is the number of proper
\(k\)-colorings. It satisfies
\[
 P_G(q)=\sum_{A\subseteq E}(-1)^{|A|}q^{c(A)},\qquad
 P_G(q)=P_{G-e}(q)-P_{G/e}(q) \quad (e\in E).
\]
For the edgeless graph on \(n\) vertices it is \(q^n\).
\end{theorem}
\begin{proof}
For each edge, declare its violation to be the event that its ends
get the same color in a function \(V\to[k]\). A function violating
all edges in \(A\) is precisely one constant on each component of
\((V,A)\). Hence there are \(k^{c(A)}\) such functions, with the
empty function counted as one. Inclusion--exclusion over the
violations gives the displayed polynomial. Two polynomials over
\(\mathbb Q\) agreeing at all nonnegative integers must be equal.
Indeed division by \(q-k\) shows that a root removes one unit of
degree, so a nonzero polynomial has at most its degree many
distinct roots.

Each proper coloring of \(G-e\) has either different colors at
\(u,v\), in which case it is a coloring of \(G\), or the same
color there. The latter colorings correspond bijectively to
colorings of \(G/e\) by descent along the identification of the
two vertices. Several copies of the same adjacency impose just
one inequality. Counting and then using polynomial uniqueness
gives deletion--contraction. If \(E\) is empty, the subset sum has
one term and \(c(\varnothing)=n\).
\end{proof}

\section{Edge colorings}

A proper edge coloring assigns different colors to distinct edges
with a common endpoint. Its minimum number of colors is the
\emph{edge chromatic number} \(\chi'(G)\). For a loopless multigraph,
parallel edges are distinct edges, and the degree counts them with
multiplicity. The edges at a vertex of maximum degree always give
the lower bound \(\chi'(G)\ge\Delta\).

\begin{theorem}[K\H{o}nig's line-coloring theorem and Vizing's bound]
\label{p204:W61}
For a finite bipartite multigraph \(\chi'(G)=\Delta(G)\).
For a finite simple graph \(\chi'(G)\le\Delta(G)+1\).
\end{theorem}
\begin{proof}
For the first claim the case \(\Delta=0\) is immediate. Pad the
two vertex parts with isolated vertices to the same size. In each
part the sum of the deficiencies \(\Delta-\deg(v)\) is then
\(\Delta\) times that common size minus the total number of edges.
While there is a deficiency on the left, there is one on the
right. Join such a pair by an edge, permitting parallel edges.
After finitely many steps the enlarged graph is \(\Delta\)-regular.
In an \(r\)-regular bipartite multigraph with \(r>0\), the \(r|S|\)
edges counted from a left subset \(S\) land in its neighborhood,
which can receive at most \(r|N(S)|\) such edges. Thus \(S\)
satisfies Hall's condition. Theorem~\ref{p204:W51}, applied to the
simple adjacency relation, gives a matching saturating the left
part. The parts have equal size, so the matching is perfect.
Choose one labelled edge for each matched pair. Removing this
perfect matching gives an \((r-1)\)-regular multigraph. Iterating
decomposes the enlarged graph into \(\Delta\) matchings, which can
be assigned separate colors. Restriction to the original graph
proves the upper bound.

For the second claim fix a palette of \(\Delta+1\) colors and
induct on the number of edges. Deleting one edge \(xv_0\) gives
a proper partial coloring with only \(xv_0\) uncolored. At
every vertex at least one color is missing. A \emph{fan} is a
list of distinct neighbors \(v_0,v_1,\ldots,v_k\) of \(x\)
such that, for \(1\le j\le k\), the color \(a_j\) of \(xv_j\)
is missing at \(v_{j-1}\). Choose a fan that cannot be extended,
which exists because the neighbor set is finite.

For a prefix through \(v_t\), rotate the colors by placing
\(a_j\) on \(xv_{j-1}\) for \(1\le j\le t\), leaving \(xv_t\)
uncolored. This is a proper partial coloring. At \(x\) the
colors on the spokes have only moved to different spokes. At
each other vertex the new color was missing there before the
rotation, and no two spokes have that vertex as an endpoint.
Here the graph being simple ensures that a chosen neighbor
determines one spoke. If some color is missing at both \(x\)
and \(v_t\) before rotating this prefix, we may color \(xv_t\)
with it after the rotation.

Choose \(\alpha\) missing at \(x\) and \(\beta\) missing at \(v_k\).
If \(\beta\) is missing at \(x\), rotate the full fan and finish.
Otherwise exactly one edge at \(x\) has color \(\beta\); its
other end must belong to the fan, since a neighbor outside
could be appended. Thus \(a_i=\beta\) for some \(1\le i\le k\).
Both \(v_{i-1}\) and \(v_k\) miss \(\beta\), and these vertices
are distinct. In the subgraph of edges with colors
\(\alpha,\beta\) each vertex has degree at most two and any
two incident edges have different colors. The component
of \(x\) is a path because \(x\) has degree one in it.
A vertex missing \(\beta\) cannot be internal to this path.
Since one endpoint is \(x\), the component cannot contain
both \(v_{i-1}\) and \(v_k\).

If \(v_{i-1}\) is outside this component, swap \(\alpha\) and
\(\beta\) on its own component, so that it misses \(\alpha\).
No edge at \(x\) changes. All spoke colors of the prefix through
\(v_{i-1}\) differ from \(\alpha\) and \(\beta\), so their
missing-color conditions persist. Rotate that prefix and color
the free edge with \(\alpha\).
If \(v_{i-1}\) is in the component of \(x\), use the component
of \(v_k\) instead. Now \(v_k\) misses \(\alpha\). The missing
\(\beta\) at \(v_{i-1}\) has not changed, and all other spoke
conditions use colors distinct from the two swapped colors.
The full fan can therefore be rotated and its last edge
colored with \(\alpha\). Swapping a whole two-color component
always preserves properness, since any incident edge of those
colors belongs to the same component. This proves the induction.
\end{proof}

For the bipartite result see \cite[Chapter 20]{Schrijver}.
The second argument uses the fan construction of \cite{MG};
the explicit choice of the component outside \(x\) keeps all
spoke colors at \(x\) fixed.

\section{Plane embeddings and five colors}

A plane embedding here draws the vertices as distinct points and
the edges as simple polygonal arcs meeting only at their prescribed
endpoints. We may regard a planar drawing as a drawing in the
sphere \(S^2\) by adjoining the point at infinity. A \emph{face}
is a connected component of the complement of the embedded graph.
In the plane the unbounded face corresponds to the component
containing the point at infinity.

\begin{theorem}[Euler's formula and planar edge bounds]\label{p204:W62}
For a finite graph embedded as above, let \(v,e,c,f\) count vertices,
edges, graph components, and faces in the sphere. Then
\[
 v-e+f=1+c.
\]
In particular, a simple planar graph with \(v\ge3\) has \(e\le3v-6\).
If it is bipartite it has \(e\le2v-4\). Thus \(K_5\) and \(K_{3,3}\)
are nonplanar.
\end{theorem}
\begin{proof}
For the empty graph \(v=e=c=0\) and \(f=1\). Otherwise its
image \(X\) is nonempty, compact, locally contractible, and proper
in \(S^2\). The local models at interior edge points and vertices
are an interval and a finite star, respectively. A finite union
of polygonal arcs and points cannot fill a two-dimensional
coordinate disk, which proves properness. Orient each edge
arbitrarily. The cellular boundary \(B:\mathbb Q^e\to\mathbb Q^v\)
sends each edge to its head minus its tail
\cite[Theorem 4.2]{P202}. The kernel of its transpose consists of
functions on vertices constant on each graph component and has
dimension \(c\). Therefore \(B\) has rank \(v-c\). There are no
two-cells, so \(H_1(X;\mathbb Q)=\ker B\) has dimension \(e-v+c\).
The universal coefficient theorem \cite[Theorem 4.13]{P202}
identifies \(H^1(X;\mathbb Q)\) with the dual of \(H_1(X;\mathbb Q)\),
since every vector space over \(\mathbb Q\) is free and the
\(\operatorname{Ext}^1_{\mathbb Q}\) term vanishes.
Alexander duality over \(\mathbb Q\)
\cite[Theorem 5.6]{P202} identifies this \(H^1\) with
\(\widetilde H_0(S^2\setminus X;\mathbb Q)\).
The complement is open and locally path connected; its path
components are its connected components. Reduced degree-zero
homology is free on the differences between those components
and a fixed one. It follows that there are finitely many faces
and \(f-1=e-v+c\).

For the inequalities first take a connected simple graph with at
least three vertices. Each edge has two sides, which may belong
to the same face. They can be traced by directed edges in the
two directions with the face kept on the same side, using the
orientation of the sphere. At a vertex continue along the first
edge in the corresponding cyclic order from the reversed
direction. The intervening sector contains no edge, so the
continuation stays on the same face. This rule is a permutation
of the directed edges, as reversing the local choice gives its
inverse. Its orbits are closed boundary walks. Every face has
at least one edge side: its boundary is nonempty and lies in
\(X\). At an interior edge point this supplies a side directly,
and at a vertex the local punctured disk minus the incident
rays has sectors bounded by edge sides. Since there are finitely
many sectors, a face approaching the vertex contains at least
one such sector. The connected graph with at least three
vertices has no isolated vertices.

A boundary walk cannot have length one in a simple graph.
If it has length two, the two directed edges must be opposite
directions of the same edge. The continuation must then reverse
immediately at both endpoints, which is possible only if both
endpoints have degree one. That would make the connected graph
a single edge. Thus every face accounts for at least three
directed sides, or at least four in a bipartite graph because
closed walks there have even length. Counting the sides gives
\(3f\le2e\) and, in the bipartite case, \(4f\le2e\).
Euler's formula with \(c=1\) yields the two bounds.

For a disconnected graph let \(a\) count isolated vertices,
\(b\) count single-edge components, and \(h\) count components
with at least three vertices. If \(h\ge1\), sum the connected
bounds. They give
\[
 e\le 3v-3a-5b-6h\le3v-6,\qquad
 e\le 2v-2a-3b-4h\le2v-4
\]
with the second estimate in the bipartite case. If \(h=0\), then
\(e=b\le\lfloor v/2\rfloor\), which implies both bounds for
\(v\ge3\). Finally \(K_5\) has ten edges on five vertices,
exceeding nine, and \(K_{3,3}\) has nine edges on six vertices,
exceeding the bipartite bound of eight.
\end{proof}

\begin{theorem}[Five-color theorem]\label{p204:W63}
Every finite simple planar graph admits a proper five-coloring.
\end{theorem}
\begin{proof}
Induct on the number of vertices, with the empty graph and the
graphs on at most two vertices immediate. For \(v\ge3\),
Theorem~\ref{p204:W62} and the degree sum show that the average
degree is less than six. Pick a vertex \(x\) of degree at most
five and color \(G-x\) by induction. If a color is unused on
its neighbors, assign that color to \(x\). In the remaining
case there are exactly five neighbors with distinct colors.
In their cyclic order about \(x\), label them \(v_1,\ldots,v_5\)
and rename the colors so that \(v_i\) has color \(i\).

If \(v_1\) and \(v_3\) lie in different components of the subgraph
induced by colors \(1,3\), interchange those two colors on the
component of \(v_1\). This frees color \(1\) at \(x\). If they
lie in the same component, take a simple path joining them
in that subgraph. This path together with \(xv_1\) and \(xv_3\)
is a simple closed polygonal curve in \(S^2\).
By Jordan separation \cite[Theorem 5.11]{P202}, it has two
complementary components with the curve as the boundary of each.
Near \(x\) its two local sectors lie in distinct components:
each sector is connected, and both global components meet
every neighborhood of \(x\). By the cyclic order, the spokes
toward \(v_2\) and \(v_4\) start in different sectors.
They do not meet the closed curve again, since the drawing
has no crossings and their other endpoints do not lie on the
\(1,3\) path. Thus \(v_2\) and \(v_4\) lie on different sides.
A path using only colors \(2,4\) in \(G-x\) would be disjoint
from that curve, since it shares neither a vertex nor a crossing
with the \(1,3\) path or its two closing edges. Such a path
cannot join \(v_2\) to \(v_4\). Interchanging colors \(2,4\)
on the component of \(v_2\) now frees color \(2\) at \(x\).

In either swap no edge with an endpoint outside the chosen
component becomes monochromatic: an endpoint outside using one
of the swapped colors would itself be in that component.
This completes the induction.
\end{proof}

For the plane graph conventions and related coloring results see
\cite[Chapters 10--11]{BM}\cite[Chapters 4--5]{Diestel}.
\chapter{Finite matroids, duality, and minors}
\label{p204:ch10}

All ground sets in these chapters are finite. We write \(A+e\) and
\(A-e\) for \(A\cup\{e\}\) and \(A\setminus\{e\}\), respectively.
All complements, unless a different ground set has been specified,
are relative to the ground set \(E\). Our starting point is an
independence system: a family \(\mathcal I\subseteq 2^E\) containing
\(\varnothing\) and closed under taking subsets. It is a
\emph{matroid} if it satisfies augmentation: for \(I,J\in\mathcal I\)
with \(|I|<|J|\), some \(e\in J\setminus I\) has \(I+e\in\mathcal I\).
A basis of a subset \(A\) means a maximal independent subset of
\(A\); a basis of the matroid is a basis of \(E\). The empty ground
set has one basis, namely \(\varnothing\).

\section{Equivalent descriptions}

\begin{theorem}[Independence, bases, rank, circuits, and closure]
\label{p204:W64}
Each of the following data on a finite \(E\) determines a matroid.
The five descriptions yield one another by the rules below.
\begin{enumerate}
\item A family \(\mathcal I\) satisfying the independence and
augmentation axioms above.
\item A nonempty family \(\mathcal B\subseteq 2^E\) with basis
exchange: for \(B,D\in\mathcal B\) and \(e\in B\setminus D\) there
is \(f\in D\setminus B\) with \(B-e+f\in\mathcal B\).
\item An integer-valued rank function \(r\) on \(2^E\) with
\(r(\varnothing)=0\), with
\(0\le r(A+e)-r(A)\le 1\), and with
\[
 r(A\cup B)+r(A\cap B)\le r(A)+r(B).
\]
\item A family \(\mathcal C\) of nonempty sets, none properly
containing another member, with circuit elimination: if \(C,D\)
are distinct members and \(e\in C\cap D\), then
\((C\cup D)-e\) contains a member of \(\mathcal C\).
\item A map \(\operatorname{cl}:2^E\to 2^E\) that is extensive,
monotone, and idempotent, with the exchange implication
\[
 y\in\operatorname{cl}(A+x)\setminus\operatorname{cl}(A)
 \quad\Longrightarrow\quad
 x\in\operatorname{cl}(A+y).
\]
\end{enumerate}
In (2) the independent sets are the subsets of members of
\(\mathcal B\). In (3) they are the sets \(I\) with \(r(I)=|I|\).
In (4) they are the sets containing no member of \(\mathcal C\).
In (5) they are the sets \(I\) with
\(e\notin\operatorname{cl}(I-e)\) for every \(e\in I\).
\end{theorem}

\begin{proof}
\emph{From independence to the other descriptions.}
Any independent subset of \(A\) extends to a basis of \(A\) by
finiteness. Augmentation forces all bases of \(A\) to have the
same size, denoted \(r(A)\). For bases \(B,D\) of \(E\) and
\(e\in B\setminus D\), apply augmentation to \(B-e\) and \(D\)
to get the required replacement from \(D\setminus B\).
Adding one element can increase rank by at most one, since
deleting it from an independent set leaves an independent set.
For submodularity, extend a basis \(I\) of \(A\cap B\) to a
basis \(J\) of \(A\), and extend \(J\) to a basis \(K\) of
\(A\cup B\). All elements of \(K\setminus J\) lie in \(B\setminus
A\), since \(J\) is already maximal in \(A\). Then
\(I\cup(K\setminus J)\) is independent in \(B\), giving
\[
 r(A\cap B)+r(A\cup B)-r(A)\le r(B).
\]
The circuits are the minimal dependent sets. For distinct
circuits \(C,D\), their intersection is independent and their
ranks are \(|C|-1,|D|-1\). Submodularity implies
\(r(C\cup D)\le |C\cup D|-2\). Thus deleting any common element
leaves a dependent set, which contains a circuit by finiteness.

Define \(\operatorname{cl}(A)=\{e\in E:r(A+e)=r(A)\}\).
Extensivity is immediate, and monotonicity follows from
submodularity, which gives, for \(A\subseteq B\),
\[
 r(B+e)-r(B)\le r(A+e)-r(A).
\]
For \(e\in B\) the left-hand side is already zero, so this
inequality also covers that case. Adding the elements of
\(\operatorname{cl}(A)\) one at a time leaves rank unchanged.
If adding \(e\) to \(\operatorname{cl}(A)\) also leaves rank
unchanged, then \(r(A+e)\le r(\operatorname{cl}(A))=r(A)\).
This proves idempotence. In the hypothesis of closure exchange,
\(r(A+y)=r(A)+1\), and \(r(A+x+y)=r(A+x)\le r(A)+1\).
Consequently \(r(A+x+y)=r(A+y)\), as required.
The stated recovery rules for \(\mathcal I\) now follow.
For the closure rule, if \(I\) is independent every deletion
drops rank by one. If \(I\) is dependent, it contains a
circuit \(C\), and for any \(e\in C\),
\(e\in\operatorname{cl}(C-e)\subseteq\operatorname{cl}(I-e)\).

\smallskip
\emph{From bases.}
First the exchange axiom forces equal sizes. Starting from
\(B\in\mathcal B\), exchange elements of \(B\setminus D\) for
elements of \(D\) until the current member is contained in \(D\).
A proper containment is impossible: exchange from \(D\) against
the smaller member would require an element outside \(D\).
Declare subsets of members independent. Suppose augmentation
fails for \(I,J\), where \(|I|<|J|\). Choose \(D\in\mathcal B\)
containing \(J\), and among members containing \(I\) choose
\(B\) maximizing \(|B\cap D|\). Then \(B\setminus D\subseteq I\),
since exchanging an element outside \(I\) improves the maximum.
Also \(J\cap B=J\cap I\), by the failure of augmentation. Hence
\[
 |D\setminus B|\ \ge\ |J\setminus B|
 > |I|-|I\cap J|\ \ge\ |I\setminus D|
 \ \ge\ |B\setminus D|.
\]
This contradicts equal sizes of \(B,D\). The resulting maximal
independent sets are exactly the original members.

\smallskip
\emph{From rank.}
The increment bounds imply \(0\le r(A)\le |A|\). A subset
\(J\subseteq I\) of an independent set is independent, since
\(r(I)\le r(J)+|I\setminus J|\). If no element of \(J\setminus I\)
increases \(r(I)\), then by the diminishing-increment inequality
above none increases rank during the successive additions, so
\(r(I\cup J)=r(I)\). This is impossible when \(r(J)=|J|>|I|\).
Thus augmentation holds. If \(I\) is maximal independent in \(A\),
no element of \(A\setminus I\) increases its rank; the same
argument gives \(r(A)=|I|\), recovering the original rank.

\smallskip
\emph{From circuits.}
Declare a set independent if it contains no designated circuit.
Whenever \(I\) is independent and \(I+x\) is dependent, there
is a unique circuit \(C_x\subseteq I+x\). Indeed, two different
circuits would both contain \(x\), and elimination would give a
circuit in \(I\). For \(y\in C_x-x\), the set \(I-y+x\) is
independent, since its only possible circuit in \(I+x\) was \(C_x\).
If \(I\) was maximal independent in \(A\) and \(x\in A\setminus I\),
then \(K=I-y+x\) is still maximal in \(A\). Adding \(y\) to \(K\)
restores \(C_x\). For \(z\in A\setminus(I+x)\), use the circuit
\(C_z\subseteq I+z\): it lies in \(K+z\) unless it contains
\(y\), and in that case elimination of \(y\) between \(C_z\)
and \(C_x\) gives a circuit in \(K+z\).

If augmentation fails for independent \(I,J\) with \(|I|<|J|\),
then \(I\) is maximal in \(A=I\cup J\). Among the maximal
independent subsets of \(A\) of size \(|I|\), choose \(I'\)
maximizing its intersection with \(J\). For \(x\in J\setminus
I'\), the circuit \(C_x\) contains a \(y\notin J\), since \(J\)
is independent. Replacing \(y\) by \(x\) gives another such
maximal independent set with larger intersection. This proves
augmentation. The antichain condition ensures that the
designated circuits are exactly the minimal dependent sets.

\smallskip
\emph{From closure.}
Independence by the closure test is hereditary by monotonicity.
If \(I\) passes the test and \(x\notin\operatorname{cl}(I)\),
then \(I+x\) passes it. A possible violation at \(x\) is excluded
by hypothesis. A violation at \(y\in I\) would put
\(y\in\operatorname{cl}((I-y)+x)\setminus\operatorname{cl}(I-y)\),
so exchange would put \(x\in\operatorname{cl}(I)\).

For an independent \(J\subseteq\operatorname{cl}(A)\) we claim
\(|J|\le |A|\). Start with the spanning set \(A\), and insert
the elements of \(J\) in sequence. At a given step let \(S\)
be the previously inserted elements and \(T\) the as yet
unremoved elements of \(A\), maintaining
\(\operatorname{cl}(S\cup T)=\operatorname{cl}(A)\). The next
element \(x\) is not in \(\operatorname{cl}(S)\) by independence.
In an enumeration of \(T\), there is a first element \(t\)
whose addition puts \(x\) in the closure. Exchange, applied to
\(S\) together with the preceding elements of the enumeration,
puts \(t\) in \(\operatorname{cl}(S\cup(T-t)+x)\).
By monotonicity and idempotence, replacing \(t\) by \(x\)
therefore preserves the closure of the union. This spends one
element of \(T\) per insertion and proves the claim.
If augmentation fails, then \(J\subseteq\operatorname{cl}(I)\),
so the claim precludes \(|J|>|I|\). Finally take a maximal
independent \(I\subseteq A\). It satisfies
\(\operatorname{cl}(I)=\operatorname{cl}(A)\), by the
one-element extension test. For \(e\in\operatorname{cl}(A)\),
the claim applied to \(I\) bounds the size of every independent
subset of \(A+e\) by \(|I|\). For
\(e\notin\operatorname{cl}(A)\), \(I+e\) is independent.
Thus the closure recovered from rank is the original one.
\end{proof}

\begin{corollary}[Bases and diminishing increments]
\label{p204:W65}
For all subsets \(A\subseteq E\), the bases of \(A\) have the
same size. The bases of \(E\) satisfy exchange. The rank is
monotone with one-element increments at most one and is
submodular. In particular an element that does not increase
rank on \(A\) does not increase it on any larger subset.
\end{corollary}

\begin{proof}
Extension and augmentation in \(A\) give equal cardinalities.
Extending \(B-e\) using a basis \(D\) gives the stated exchange.
The extension from \(A\cap B\) through \(A\) to \(A\cup B\) in
Theorem~\ref{p204:W64} proves the rank assertions; its
diminishing-increment inequality gives the last claim.
\end{proof}

\begin{corollary}[Circuit elimination and rank closure]
\label{p204:W66}
The circuits are precisely the inclusion-minimal dependent sets.
Two distinct circuits sharing \(e\) have a circuit in their
union with \(e\) removed. For \(e\in E\), \(e\) belongs to the
closure of \(A\) exactly when adding it leaves the rank unchanged.
\end{corollary}

\begin{proof}
These are the circuit and closure constructions in
Theorem~\ref{p204:W64}. In particular, the rank of the circuit
union is at most its size minus two, even when the circuits
intersect in more than one element. Removing the specified
element therefore gives a dependent set.
\end{proof}

\section{Duality and elementary structure}

\begin{theorem}[Duality, deletion, and contraction]
\label{p204:W67}
The complements \(E\setminus B\) of the bases of \(M\) are the
bases of a matroid \(M^*\) on \(E\). Duality is an involution and
\[
 r_{M^*}(A)=|A|+r_M(E\setminus A)-r_M(E)
 \qquad (A\subseteq E).
\]
For \(C\subseteq E\), define deletion \(M\setminus C\) by
retaining the independent sets contained in \(E\setminus C\),
and contraction \(M/C=(M^*\setminus C)^*\) on \(E\setminus C\).
Then
\[
 (M/C)^*=M^*\setminus C,\qquad
 (M\setminus C)^*=M^*/C.
\]
\end{theorem}

\begin{proof}
For two primal bases \(B,D\), let \(e\in D\setminus B\).
There is a unique circuit \(C_e\) in \(B+e\) by the elimination
argument in Theorem~\ref{p204:W64}. It contains \(e\) and an
\(f\in B\setminus D\), since \(D\) is independent. Then
\(B-f+e\) is independent and has full rank, hence is a basis.
Its complement is the result of exchanging \(e\) for \(f\)
in \(E\setminus B\), proving the dual basis axiom. Taking
complements twice gives the original bases.

Every independent set extends to a basis, so
\(r_M(A)=\max_B|B\cap A|\). Therefore
\(\min_B|B\cap A|=r_M(E)-r_M(E\setminus A)\), and maximizing
\(|(E\setminus B)\cap A|=|A|-|B\cap A|\) gives the rank formula.
Deletion preserves the axioms because both the heredity test
and augmentation use elements only from the retained sets.
The two displayed duality identities now follow from the
definition and involutivity on the smaller ground set.
\end{proof}

\begin{theorem}[Fundamental circuits, cocircuits, and simplification]
\label{p204:W68}
Let \(I\) be independent and \(e\in\operatorname{cl}(I)\setminus I\).
There is exactly one circuit in \(I+e\), its \emph{fundamental
circuit}. For \(f\in I\), the set \(I-f+e\) is independent
exactly when \(f\) lies in that circuit.

A loop is an element whose singleton is a circuit, and a
coloop is an element belonging to every basis. A cocircuit,
defined as a circuit of \(M^*\), is a minimal set meeting
every basis of \(M\). In particular the coloops are the loops
of \(M^*\). Among nonloops, define two distinct elements to
be parallel if their pair has rank one. Equality or parallelism
is an equivalence relation. Choosing one representative of
each class and deleting all other elements gives a simple
matroid, meaning that it has neither loops nor parallel
distinct elements. If \(S\) is the representative set and
\(q(A)\subseteq S\) is the set of representatives of the
nonloops in \(A\), then
\[
 r_{M|S}(q(A))=r_M(A),\qquad
 \operatorname{cl}_{M|S}(q(A))=\operatorname{cl}_M(A)\cap S .
\]
\end{theorem}

\begin{proof}
Here \(I+e\) is dependent by rank closure. Any circuit in it
must contain \(e\). Elimination between two such circuits
would give a circuit in \(I\). Deleting \(f\) removes the
unique obstruction exactly when \(f\) lies in it.

A subset is independent in \(M^*\) precisely when it lies in
a complementary basis, that is, when it misses some primal
basis. Minimal dual dependence is thus minimal intersection
with every primal basis. The assertion about coloops is the
singleton case. Equivalently, a cocircuit is the complement
of a maximal nonspanning subset: meeting every basis says
exactly that the complement contains no basis, and the two
extremal conditions reverse under complementation.

For nonloops, a pair has rank one exactly when both elements
have the same singleton closure. One direction follows
because each then belongs to the closure of the other, and
the reverse follows from the rank test. This proves
transitivity. Loops are already in the closure of
\(\varnothing\). Thus every element of \(A\) lies in the
closure of \(q(A)\), and every representative in \(q(A)\)
lies in the closure of \(A\). Monotonicity and idempotence
give equality of the two closures and of the ranks.
Restriction to \(S\) preserves the rank of every subset of
\(S\), directly by the independent-set definition. Testing
the rank after adding one element of \(S\) yields the closure
formula. A singleton in \(S\) has rank one and two distinct
members cannot be parallel, proving simplicity.
\end{proof}

\section{Rank under minors}

\begin{theorem}[Rank and operations on minors]
\label{p204:W69}
For \(A\subseteq E\setminus C\),
\[
 r_{M\setminus C}(A)=r_M(A),\qquad
 r_{M/C}(A)=r_M(A\cup C)-r_M(C).
\]
If \(I\) is any basis of \(C\), then \(J\subseteq E\setminus C\)
is independent in \(M/C\) exactly when \(I\cup J\) is independent
in \(M\). Successive deletions combine, successive contractions
combine, and deletion and contraction on disjoint sets commute.
These assertions include empty deleted or contracted sets and
the cases in which the remaining ground set is empty.
\end{theorem}

\begin{proof}
Deletion gives the first formula by inspection of independent
sets. To compute contraction rank, put \(E'=E\setminus C\).
Applying dual rank on \(E'\), deletion rank, and then dual
rank on \(E\), in that order, gives
\begin{align*}
 r_{M/C}(A)
 &= |A|+r_{M^*}(E'\setminus A)-r_{M^*}(E')\\
 &= |A|+|E'\setminus A|-|E'|
       +r_M(C\cup A)-r_M(C)\\
 &=r_M(C\cup A)-r_M(C).
\end{align*}
The set \(C\) lies in the closure of \(I\); adding \(J\) to
both sets therefore leaves
\(r_M(I\cup J)=r_M(C\cup J)\). Since \(I\) is disjoint from
\(J\) and \(|I|=r_M(C)\), the rank test for \(I\cup J\) is
exactly the contraction rank test for \(J\). This also
explains the formula as extension of a basis of \(C\).

Successive deletions give restriction to the final ground
set. For disjoint sets \(C,D\), applying the contraction
formula twice gives
\[
 r_{(M/C)/D}(A)=r_M(A\cup C\cup D)-r_M(C\cup D).
\]
With one deletion and one contraction on disjoint sets, the
rank on the remaining ground set is
\(r_M(A\cup C)-r_M(C)\) in either order, where \(C\) is
contracted. Equality of rank functions on the same ground
set gives equality of matroids by Theorem~\ref{p204:W64}.
The formulas use the zero rank of the empty set and
require no rank-positivity assumptions.
\end{proof}

For comparison of the standard presentations and conventions,
see \cite[Sections 1.1--1.4 and Chapters 2--3]{Oxley}.
\chapter{Representations, finite geometry, and enumeration}
\label{p204:ch11}

A representation records dependence among \emph{indexed} columns.
A zero column, repeated columns, and distinct proportional columns
can therefore carry different index information. In this chapter
a field need not have characteristic zero.

\section{Vector configurations and graphical representations}

\begin{theorem}[The matroid of a vector family]
\label{p204:W70}
Let \(K\) be a field, \(V\) a \(K\)-vector space, and
\((v_e)_{e\in E}\) a finitely indexed family in \(V\).
The sets \(I\subseteq E\) on which the columns are linearly
independent form a matroid. Its rank on \(A\subseteq E\) is
\[
 r(A)=\dim_K\operatorname{span}\{v_e:e\in A\}.
\]
\end{theorem}

\begin{proof}
Heredity is immediate. For independent \(I,J\) with \(|I|<|J|\),
the span of the \(J\)-columns cannot be contained in the span
of the \(I\)-columns, by the dimension inequality for nested
finite-dimensional spaces. A \(J\)-column outside the latter
span augments \(I\). A maximal independent subfamily of \(A\)
spans all its columns, since any column outside its span could
be added. Its cardinality is thus the displayed dimension.
The argument applies to repeated and zero columns. In particular,
zero columns are loops.
\end{proof}

\begin{theorem}[Deletion, contraction, and extension of scalars]
\label{p204:W71}
Let the matroid of \((v_e)_{e\in E}\) over \(K\) be \(M\).
Deleting \(C\subseteq E\) is represented by omitting the
\(C\)-columns. Contracting \(C\) is represented by the images
of the remaining columns in
\[
 V/\operatorname{span}\{v_e:e\in C\}.
\]
For any field embedding \(K\hookrightarrow L\), extension
of scalars from \(K\) to \(L\) preserves the rank of every
subfamily, hence preserves the matroid.
\end{theorem}

\begin{proof}
Deleting columns gives exactly the same independence tests
on the smaller ground set. Set \(U=\operatorname{span}\{v_e:e\in
C\}\) and \(W=\operatorname{span}\{v_e:e\in A\}\) for
\(A\subseteq E\setminus C\). The images of the \(A\)-columns
span \((U+W)/U\). A basis of \(U\) extends to a basis of
\(U+W\); the newly added vectors project to a basis of this
quotient, by lifting a linear relation and using the
independence of the extended basis. Its dimension is
\(\dim(U+W)-\dim U=r_M(A\cup C)-r_M(C)\).
Theorem~\ref{p204:W69} identifies that rank as contraction rank.

For scalar extension, use coordinates in a finite basis of the
span of the entire family. This finite span has a linear
retraction from \(V\), obtained by extending its basis to one of
\(V\) and projecting. Thus its inclusion remains injective after
extension of scalars, since the same retraction extends. In the
chosen coordinates, ordinary elimination over \(K\) brings any
column submatrix to a pivot form by invertible row and column
operations. These operations remain invertible over \(L\),
and the pivot form has the same rank there. This proves rank
preservation, also when the number of pivots is zero.
The corresponding finite-matrix scalar extension and quotient
dimension statements are available in the D05 development
\cite{D05}.
\end{proof}

For a finite undirected graph, its components include isolated
vertices. A loop is regarded as a cycle of length one, and a pair
of distinct parallel edges as a cycle of length two. This
convention extends the simple graphs used in Chapter~\ref{p204:ch06}.
A forest contains no cycle of any of these kinds. For a given
orientation of each nonloop edge, the incidence column of
\(u\longrightarrow v\) is \(\mathbf e_v-\mathbf e_u\) in the
vertex coordinate space. Give each loop the zero column.

A \emph{flat} of a matroid is a set \(F\) with
\(\operatorname{cl}(F)=F\). In a lattice \(L\), write \(x\prec y\)
when \(y\) covers \(x\), meaning \(x<y\) with no element
strictly between. Call a finite lattice \emph{geometric} here
if each element is a join of atoms (elements covering \(0\)),
and it is upper semimodular in the following covering form:
if \(x\) covers \(x\wedge y\), then \(x\vee y\) covers \(y\).
The empty join is \(0\).

\begin{theorem}[Cycle matroids and geometric lattices]
\label{p204:W72}
For a finite graph, over every field and every choice of
orientation, the incidence columns represent the matroid
whose independent sets are the forests. Its bases are the
unions of spanning trees of the components, and its rank is
\(v-c\), where \(v\) is the number of vertices and \(c\) the
number of components.

The flats of every finite matroid form a geometric lattice.
Conversely, the atoms of a finite geometric lattice carry a
simple matroid whose flat lattice is the given lattice.
This construction of the simple matroid uses only the lattice.
\end{theorem}

\begin{proof}
A cycle yields a linear dependence by following it and summing
its oriented columns, giving each column coefficient \(1\) or
\(-1\) according to the traversal direction. The sum telescopes
to zero and every coefficient on the cycle is nonzero, even
in characteristic two. For a loop the zero column gives the
dependence directly. Conversely, a nonempty forest with at
least one edge has a leaf in an edge-containing component:
an endpoint of a longest simple path cannot have an additional
neighbor, because that would extend the path or create a cycle.
At the leaf, a linear relation forces the coefficient of its
unique incident edge to vanish. Removing that edge and
inducting proves independence.

An edge joining two components of a forest may be added
without creating a cycle. Thus a maximal forest must connect
all vertices of each original connected component, or an edge
of the graph between two of its forest components would still
be available. Conversely, adding an edge within one tree
creates a cycle. A tree on \(n\geq 1\) vertices has \(n-1\)
edges, by deleting a leaf and inducting, with the singleton
as the base case. Summing over all components gives \(v-c\).
This includes the empty graph and the edgeless graph.

For the flat lattice, intersections are meets and
\(F\vee G=\operatorname{cl}(F\cup G)\) is the join. The bottom
is \(\operatorname{cl}(\varnothing)\). If \(F\subsetneq G\)
are flats and \(x\in G\setminus F\), then
\(\operatorname{cl}(F+x)\) is between them, strictly above
\(F\), with rank \(r(F)+1\). It follows that covers are
exactly strict flat inclusions with rank difference one.
The atoms are precisely the closures of nonloop singletons,
and every flat is their join over the nonloops it contains.
For the semimodular condition, if \(F\) covers \(F\cap G\),
rank submodularity gives
\[
 r(F\vee G)\le r(F)+r(G)-r(F\cap G)=r(G)+1.
\]
The join is strictly above \(G\), since \(F\) is not contained
in \(G\). Hence it covers \(G\).

For the converse, first establish the rank of the lattice
itself. If \(u\prec v\), then for any \(b\in L\) either
\(u\vee b=v\vee b\) or \(u\vee b\prec v\vee b\). In the
unequal case \(v\wedge(u\vee b)=u\): this meet lies between
\(u,v\) and cannot equal \(v\). Upper semimodularity then
gives the lifted cover. Every interval has maximal chains
of equal length, by induction on its size. Indeed, for two
chains starting at the same lower endpoint \(a\), if their
first covers agree, use the shorter interval above the
common cover. If the first covers \(p,q\) differ, then
\(p\wedge q=a\) and \(p\vee q\) covers both. Extend a chain
from \(p\vee q\) to the upper endpoint. Induction in the
intervals starting at \(p\) and at \(q\) compares the two
original chain lengths to the lengths of the chains through
this common join, which agree. The one-element interval has
length zero.

Write \(h(x)\) for the length of such a chain from \(0\)
to \(x\). For \(x\le y\) the length in \([x,y]\) is
\(h(y)-h(x)\) by concatenation. Lift a saturated chain from
\(x\wedge y\) to \(x\) by joining with \(y\) at each step.
Each step collapses or is a cover. The resulting chain
therefore gives
\[
 h(x\vee y)-h(y)\le h(x)-h(x\wedge y).
\]
On the set of atoms define \(r(A)=h(\bigvee A)\). The lifted
cover property for \(0\prec a\) gives the one-atom increment
bound. The lattice inequality gives submodularity, using
\(\bigvee(A\cap B)\le(\bigvee A)\wedge(\bigvee B)\).
Thus this is a matroid rank function by
Theorem~\ref{p204:W64}. A single atom has rank one and
distinct atoms have join of rank two, so the matroid is
simple. For an atom \(a\), adjoining it preserves rank
exactly when \(a\le\bigvee A\), since \(h\) strictly
increases along proper inclusions. Hence the flats are the
sets of atoms below lattice elements. Taking joins and
taking atoms below give inverse order-preserving maps
because the lattice is generated by its atoms.
\end{proof}

\section{Field restrictions and projective examples}

\begin{theorem}[A rank-two obstruction]
\label{p204:W73}
The uniform rank-two matroid \(U_{2,4}\), with every set of
size at most two independent, is representable over \(K\)
exactly when \(K\) has at least three elements. More generally,
a simple rank-two matroid represented over \(\mathbb F_q\)
has at most \(q+1\) elements.
\end{theorem}

\begin{proof}
Replace the ambient space by the two-dimensional column span.
Simplicity requires every column to be nonzero and any two
columns to generate distinct one-dimensional subspaces.
Over \(\mathbb F_q\), the nonzero vectors are partitioned by
these subspaces, each of which contributes \(q-1\) nonzero
vectors. Thus there are
\((q^2-1)/(q-1)=q+1\) possible subspaces. In particular
\(\mathbb F_2\) has only three. Conversely, over any field
with \(t\notin\{0,1\}\), the four columns
\[
 (1,0)^{\mathsf T},\quad(0,1)^{\mathsf T},\quad
 (1,1)^{\mathsf T},\quad(1,t)^{\mathsf T}
\]
determine four distinct one-dimensional subspaces, which
represent \(U_{2,4}\).
\end{proof}

\begin{theorem}[The Fano characteristic criterion]
\label{p204:W74}
The Fano matroid, represented by the seven nonzero vectors
of \(\mathbb F_2^3\), is representable over a field \(K\)
if and only if \(\operatorname{char}K=2\).
\end{theorem}

\begin{proof}
When \(K\) has characteristic two, the copy of \(\mathbb F_2\)
generated by \(1\) embeds in \(K\), and
Theorem~\ref{p204:W71} proves sufficiency. Conversely label the
Fano points
\[
 a_1=e_1,\quad a_2=e_2,\quad a_3=e_3,\quad
 d=e_1+e_2,\quad f=e_1+e_3,\quad
 g=e_2+e_3,\quad h=e_1+e_2+e_3.
\]
In any representation the \(a_i\) are a basis of the column
span. Make them the coordinate axes. The triples
\(\{a_1,a_2,d\}\), \(\{a_1,a_3,f\}\),
and \(\{a_2,a_3,g\}\) are circuits, so each displayed
two-axis column has both indicated coordinates nonzero.
By scaling the second and third coordinates, then rescaling
columns, arrange
\[
 d=(1,1,0)^{\mathsf T},\qquad f=(1,0,1)^{\mathsf T},
 \qquad g=(0,1,t)^{\mathsf T},\quad t\ne0 .
\]
The basis columns may be rescaled to restore unit coordinates;
all these operations preserve dependence. The circuits
\(\{d,a_3,h\}\) and \(\{f,a_2,h\}\) force the first and
second, and the first and third, coordinates of \(h\)
to agree. Because \(h\) is nonzero it can be scaled to
\((1,1,1)^{\mathsf T}\). The circuit \(\{g,a_1,h\}\) now
forces \(t=1\). Finally \(\{d,f,g\}\) is a circuit in the
Fano matroid, so the determinant of these three normalized
columns must vanish. It is \(-2\), which forces
\(\operatorname{char}K=2\).
\end{proof}

\begin{theorem}[Dual columns and finite geometries]
\label{p204:W75}
If a rank-\(r\) matroid is represented by a matrix
\([I_r\ \ A]\) with \(n\) columns, its dual is represented,
with the same column indexing, by
\([-A^{\mathsf T}\ \ I_{n-r}]\).

For \(r\geq 1\), the projective points over \(\mathbb F_q\) in
dimension \(r-1\), represented by a nonzero vector from each
one-dimensional subspace of \(\mathbb F_q^r\), form a simple
rank-\(r\) matroid on \((q^r-1)/(q-1)\) points.
A projective line, when \(r\ge2\), has \(q+1\) points.
The affine points of \(\mathbb F_q^d\), represented by
\((1,x)\in\mathbb F_q^{d+1}\), form a simple rank-\((d+1)\)
matroid on \(q^d\) points. For \(d\ge1\), an affine line
has \(q\) points.
\end{theorem}

\begin{proof}
The kernel of the linear map with matrix \([I_r\ \ A]\)
consists of vectors \((-Ay,y)\). Thus the rows of
\([-A^{\mathsf T}\ \ I_{n-r}]\) form a basis of that kernel
as coordinate vectors. If \(S\) is a set of \(r\) columns,
the original columns on \(S\) are independent exactly when
no nonzero kernel vector is supported on \(S\). This is
equivalent to injectivity of projection of the kernel onto
the complementary \(n-r\) coordinates, and hence to that
projection being an isomorphism by equal dimensions. Its
matrix in the indicated kernel basis is the transpose of
the proposed dual submatrix on the complementary columns.
Thus the bases of the proposed dual are exactly the
complements of the original bases. This argument also
works for \(r=0\) or \(r=n\).

For projective points, the count partitions the \(q^r-1\)
nonzero vectors into classes of size \(q-1\). Distinct
classes give nonproportional columns and the coordinate
axis classes already span the space. A projective line
consists of the classes in a two-dimensional subspace,
so the same count gives \(q+1\). For affine points,
\((1,x)\) and \((1,y)\) are proportional only when
\(x=y\). The columns for \(0\) and the \(d\) coordinate
vectors span \(\mathbb F_q^{d+1}\), since subtracting the
column for \(0\) gives the other coordinate directions.
When \(d=0\) the single column already spans. An affine
line is \(\{x+ta:t\in\mathbb F_q\}\) for \(a\ne0\), which
has \(q\) distinct elements. The normalized columns on
such a line lie in the span of \((1,x),(0,a)\) and are
exactly its points with first coordinate normalized to one.
\end{proof}

\section{Matching and enumeration bridges}

\begin{theorem}[Transversal matroids]
\label{p204:W76}
Let \(E\) and \(R\) be the sides of a finite bipartite
relation. The sets \(I\subseteq E\) that can be saturated
by a matching into \(R\) form a matroid on \(E\).
\end{theorem}

\begin{proof}
The empty set is matchable and restriction of a matching
gives heredity. For two matchable sets \(I,J\) with
\(|I|<|J|\), choose matchings \(M_I,M_J\) whose left supports
are exactly \(I,J\), discarding edges outside those
supports if necessary. In the symmetric difference,
each vertex has degree at most two, and successive edges
alternate between the matchings. Consequently each
nontrivial component is an alternating path or an even
alternating cycle. The total number of \(M_J\)-edges is
larger than the total number of \(M_I\)-edges, so some path
has one more \(M_J\)-edge than \(M_I\)-edge. Both its end
edges belong to \(M_J\). The path has odd length and hence
one endpoint on each side. Its left endpoint \(x\) is
in \(J\setminus I\), and its right endpoint is unmatched
by \(M_I\). Replacing the \(M_I\)-edges of this path by
the \(M_J\)-edges leaves every internally matched vertex
matched and adds the two endpoints. The resulting matching
saturates \(I+x\). This proves augmentation.
\end{proof}

\begin{theorem}[The Tutte polynomial]
\label{p204:W77}
For a finite matroid \(M\) on \(E\), let
\[
 T_M(x,y)=\sum_{A\subseteq E}
  (x-1)^{r_M(E)-r_M(A)}(y-1)^{|A|-r_M(A)} .
\]
For \(e\in E\) the following identities hold:
\[
 T_M=
 \begin{cases}
   yT_{M\setminus e},& e\text{ is a loop},\\
   xT_{M/e},& e\text{ is a coloop},\\
   T_{M\setminus e}+T_{M/e},& e\text{ is neither}.
 \end{cases}
\]
If \(M_1\oplus M_2\) is the direct sum on disjoint ground
sets, then \(T_{M_1\oplus M_2}=T_{M_1}T_{M_2}\).
The evaluation \(T_M(1,1)\) is the number of bases.

For a finite simple graph \(G\), let \(M(G)\) be its cycle
matroid, let \(c(G)\) count all components, and let
\(\chi_G(t)\) be its chromatic polynomial. Then
\[
 \chi_G(t)
   =(-1)^{r_{M(G)}(E)}\,t^{c(G)}T_{M(G)}(1-t,0).
\]
\end{theorem}

\begin{proof}
The exponents are nonnegative by monotonicity and
\(r_M(A)\le|A|\). Write \(E'=E-e\), and group the summands
for \(A\) and \(A+e\) with \(A\subseteq E'\). If \(e\) is
neither a loop nor a coloop, \(r_M(\{e\})=1\) and
\(r_M(E')=r_M(E)\). The absent-\(e\) terms are those of
\(T_{M\setminus e}\). For a present-\(e\) term,
Theorem~\ref{p204:W69} gives
\[
 \begin{split}
 r_{M/e}(E')-r_{M/e}(A)&=r_M(E)-r_M(A+e),\\
 |A|-r_{M/e}(A)&=|A|+1-r_M(A+e),
 \end{split}
\]
so these terms give \(T_{M/e}\).

If \(e\) is a loop, it never increases rank, by diminishing
increments from the empty set. The two terms then differ by
a factor \(y-1\), giving the factor \(1+(y-1)=y\).
If \(e\) is a coloop, deleting it drops total rank by one:
deleting it from a basis gives the lower bound, and a
full-rank subset avoiding \(e\) would contain a basis
avoiding it. Thus \(e\) is outside the closure of \(E'\),
and outside the closure of every \(A\subseteq E'\).
It always increases rank by one on these \(A\). Contraction
and deletion have the same rank function on \(E'\).
The absent and present terms then give respectively the
factors \(x-1\) and \(1\) times the contraction term.
This proves the coloop formula.

In a direct sum, independence is tested separately on the
two disjoint ground sets. The augmentation axiom follows
by augmenting in a component in which the second
independent set has larger size. Ranks add, so splitting
\(A\) into its two components factors the defining sum.
At \(x=y=1\), a term is one exactly when both exponents
are zero. That means \(A\) is independent and has full
rank, exactly the basis condition. The rank-zero case and
the empty ground set use the constant monomial \(1\).

For the graph specialization, put \(v=|V(G)|\),
\(r=r_{M(G)}(E)=v-c(G)\), and write \(c(A)\) for the
number of components of the spanning graph with edge set
\(A\). By Theorem~\ref{p204:W72}, \(r_{M(G)}(A)=v-c(A)\).
Substituting \(x=1-t,y=0\) in the subset sum, a term
acquires sign \((-1)^{r-r(A)+|A|-r(A)}=(-1)^{r+|A|}\).
After multiplying by \((-1)^r t^{c(G)}\) it becomes
\((-1)^{|A|}t^{c(A)}\). This is the edge-subset
inclusion--exclusion formula of Theorem~\ref{p204:W60}:
requiring the endpoints of every edge in \(A\) to have
equal colors allows one free color per component.
Thus the polynomial identity follows, including for
the empty graph.
\end{proof}

The representation and geometry conventions may be compared with
\cite[Sections 1.5--1.7 and Chapters 5--6, 9]{Oxley}. The
rank-polynomial bridge belongs to the framework of \cite{Tutte}.
\chapter{Weighted greedy selection and its converse}
\label{p204:ch12}

For a weight function \(w:E\to\mathbb R\) put
\(w(A)=\sum_{e\in A}w(e)\). In an independence system a
\emph{maximal} independent set means one maximal by inclusion.
Its cardinality need not be the maximum possible cardinality
unless the augmentation axiom holds.

\begin{definition}[Weighted greedy selection]
\label{p204:W78}
Fix a total tie order on \(E\), and list the elements as
\(e_1,\dots,e_n\) in nonincreasing weight, using that order
when the weights agree. Put \(I_0=\varnothing\). At step
\(i\), let \(I_i=I_{i-1}+e_i\) if this set is independent,
and \(I_i=I_{i-1}\) otherwise. Return \(I_n\).
The \emph{basis objective} maximizes weight over maximal
independent sets. The \emph{independent-set objective}
maximizes weight over all independent sets.
\end{definition}

\begin{lemma}[The prefix invariant]
At each step \(i\), the selected set \(I_i\) is maximal
independent in the prefix \(P_i=\{e_1,\dots,e_i\}\).
If the independence system is a matroid, then
\[
 |I_n\cap P_i|=|I_i|=r(P_i).
\]
\end{lemma}

\begin{proof}
Every selected set is independent by its construction.
If \(e_j\in P_i\setminus I_i\), then \(e_j\) was rejected
at step \(j\), so \(I_{j-1}+e_j\) was dependent.
Since \(I_{j-1}\subseteq I_i\), heredity excludes
independence of \(I_i+e_j\). Thus \(I_i\) is maximal in
the prefix. Later steps do not add anything to that prefix.
In a matroid every maximal independent set of \(P_i\)
has size \(r(P_i)\) by Theorem~\ref{p204:W64}.
\end{proof}

\begin{theorem}[The matroid greedy theorem]
\label{p204:W79}
For every real weight function on a finite matroid, the
algorithm of Definition~\ref{p204:W78} returns a
maximum-weight basis, with any tie order.
If all weights are nonnegative, its output also maximizes
weight over independent sets. For arbitrary real weights,
to maximize weight over independent sets, first delete all
negative-weight elements and apply greedy to the restriction.
\end{theorem}

\begin{proof}
If \(E\) is empty, its only basis and only independent
set is the empty set. Otherwise write \(w_i=w(e_i)\).
For any \(B\subseteq E\), summing the increments of
\(|B\cap P_i|\) gives
\[
 w(B)
 = w_n|B|+\sum_{i=1}^{n-1}
     (w_i-w_{i+1})\,|B\cap P_i| .
\]
Indeed an element in position \(j\) contributes
\(w_n+\sum_{i=j}^{n-1}(w_i-w_{i+1})=w_j\).
All differences in the sum are nonnegative.
For a basis \(B\), \(|B\cap P_i|\le r(P_i)\), and the
prefix invariant attains equality for the greedy basis
in every prefix. The last term \(w_n|B|\) is constant
among bases since their sizes agree; its sign is
irrelevant. Hence greedy attains at least the weight
of every basis.

For nonnegative weights an independent set may be
extended to a basis without decreasing weight, so the
same basis also optimizes over all independent sets.
For arbitrary weights, deleting negative-weight elements
from any independent set increases its weight and leaves
it independent. Thus the optimum is achieved in the
restriction to nonnegative elements, where the preceding
argument applies. For example, a one-element coloop with
weight \(-1\) is required in a basis but is omitted by the
optimal independent set.
\end{proof}

\begin{theorem}[Greedy characterizes matroid augmentation]
\label{p204:W80}
Let \(\mathcal I\) be a nonempty hereditary family on a
finite ground set. The following are equivalent.
\begin{enumerate}
\item \(\mathcal I\) satisfies augmentation.
\item For every nonnegative weight function, weighted
greedy maximizes weight over maximal independent sets.
\item For every nonnegative weight function, weighted
greedy maximizes weight over all independent sets.
\end{enumerate}
In (2) and (3), the equivalence holds whether one fixes a
tie order in advance or requires optimality for every tie
order.
\end{theorem}

\begin{proof}
A nonempty hereditary family contains \(\varnothing\).
For nonnegative weights any independent set extends to
a maximal independent set without loss of weight.
The greedy output is maximal by the prefix invariant, so
the two optimality statements are equivalent.
Augmentation implies both by Theorem~\ref{p204:W79}.

For the converse suppose that \(I,J\in\mathcal I\) have
\(|I|<|J|\) but no element of \(J\setminus I\) augments
\(I\). Choose \(\varepsilon=1/(|I|+1)\) and put
\[
 w(e)=
 \begin{cases}
  1+\varepsilon,& e\in I,\\
  1,& e\in J\setminus I,\\
  0,& \text{otherwise}.
 \end{cases}
\]
Greedy first selects every element of \(I\), because each
prefix of \(I\) is independent. It then rejects every
element of \(J\setminus I\) by the failed augmentation.
Further selections have weight zero. The returned weight
is \((1+\varepsilon)|I|<|I|+1\le |J|\le w(J)\),
so greedy is suboptimal. The strict priority of the
three weight levels makes this argument independent of
tie order. Thus optimality forces augmentation.
\end{proof}

The finite exchange and prefix-rank arguments explain why
the basis and independent-set formulations must specify
different conventions for negative weights.
See \cite[Section 1.8]{Oxley} and \cite[Chapter 40]{Schrijver}
for related formulations.

\backmatter
\begin{thebibliography}{99}
\bibitem{Stanley} R. P. Stanley, \emph{Enumerative Combinatorics},
  vol.~1, 2nd ed., Cambridge University Press, 2012. The archive
  contains the author preprint dated 21 February 2011.
\bibitem{GKP} R. L. Graham, D. E. Knuth, and O. Patashnik,
  \emph{Concrete Mathematics}, 2nd ed., Addison--Wesley, 1994.
\bibitem{FS} P. Flajolet and R. Sedgewick, \emph{Analytic Combinatorics},
  Cambridge University Press, 2009.
\bibitem{Polya} G. P\'olya and R. C. Read,
  \emph{Combinatorial Enumeration of Groups, Graphs, and Chemical Compounds},
  Springer, 1987.
\bibitem{LGV} B. Lindstr\"om, On the vector representations of induced
  matroids, \emph{Bull. London Math. Soc.} 5 (1973), 85--90;
  I. Gessel and G. Viennot, Binomial determinants, paths, and hook length
  formulae, \emph{Adv. Math.} 58 (1985), 300--321.
\bibitem{Postnikov} A. Postnikov, \emph{18.212 Algebraic Combinatorics},
  MIT course materials, Spring 2025.
\bibitem{BM} J. A. Bondy and U. S. R. Murty, \emph{Graph Theory},
  Graduate Texts in Mathematics 244, Springer, 2008.
\bibitem{Diestel} R. Diestel, \emph{Graph Theory}, 5th ed., Springer, 2017.
\bibitem{Schrijver} A. Schrijver, \emph{Combinatorial Optimization:
  Polyhedra and Efficiency}, Springer, 2003.
\bibitem{MG} J. Misra and D. Gries, A constructive proof of Vizing's
  theorem, \emph{Inform. Process. Lett.} 41 (1992), 131--133.
\bibitem{Oxley} J. Oxley, \emph{Matroid Theory}, 2nd ed.,
  Oxford University Press, 2011.
\bibitem{Tutte} W. T. Tutte, A contribution to the theory of chromatic
  polynomials, \emph{Canad. J. Math.} 6 (1954), 80--91.
\bibitem{D05} MathGenome, \emph{Unified Linear Algebra Notes},
  D05, released with MathGenome v0.8.0.
\bibitem{D07} MathGenome, \emph{Unified Abstract Algebra Notes},
  D07, released with MathGenome v0.8.0.
\bibitem{P202} MathGenome, \emph{Foundations of Algebraic Topology
  and Riemannian Geometry}, P2-02 unified notes, 8 November 2026.
\bibitem{P203} MathGenome, \emph{Foundations of Elementary Number Theory},
  P2-03 unified notes, 8 November 2026.
\bibitem{Mathlib} The Mathlib Community, \emph{Mathlib}, pinned Lean library.
  The precise source revision is recorded in the result map.
\end{thebibliography}
\end{document}
